% =====================================================================
%  Parte 2 — Programación dinámica
%  Fragmento: se incluye con \input desde el documento maestro.
%  Macros locales con prefijo pd.
% =====================================================================

\newcommand{\pdalgo}[1]{\par\medskip\noindent{\sans\small\bfseries\color{acento}#1}\par\nopagebreak\smallskip}
\newcommand{\pdargmin}{\operatorname*{arg\,min}}
\newcommand{\pdargmax}{\operatorname*{arg\,max}}
\newcommand{\pdbar}[1]{\overline{#1}}

\parte{pd}{Programación dinámica}{\OPT(j)=\max\{v_j+\OPT(p(j)),\ \OPT(j-1)\}}{Se arma una familia chica de subproblemas ordenados, se escribe cómo sale el óptimo de cada uno a partir de los más chicos y se llena una tabla. Es el tema que más se pregunta: hubo un ejercicio de programación dinámica en todos los segundos parciales desde 2018.}

% =====================================================================
\section{El método y cómo se corrige}
% =====================================================================

La programación dinámica (PD) sirve cuando el problema se puede escribir en
función de una cantidad polinomial de subproblemas que cumplen tres cosas:
tienen un orden natural de ``más chico'' a ``más grande''; el óptimo de cada
uno se obtiene con una cuenta sencilla a partir de óptimos de subproblemas
más chicos; y el problema original es uno de ellos (o sale de ellos de forma
directa). A diferencia de dividir y conquistar, los subproblemas se
\emph{solapan}: el mismo subproblema aparece en muchas ramas de la recursión,
así que se calcula una sola vez y se guarda en una tabla.

La propiedad que lo hace funcionar es el \emph{principio de optimalidad}: una
solución óptima de una instancia está formada por soluciones óptimas de las
subinstancias en que se descompone. Se demuestra siempre igual, ``cortando y
pegando'': si la parte que corresponde a la subinstancia no fuera óptima, la
reemplazo por una mejor y obtengo una solución mejor de la instancia
original, lo cual es absurdo.\visto{Examen feb. 2013 Ej. 2: enunciar principio de optimalidad}

En los exámenes viejos (hasta 2017) se pedía enunciarlo y demostrarlo para
un problema dado, por ejemplo la mochila o los caminos mínimos de Floyd. En los
actuales no se pide con ese nombre, pero la justificación de la recurrencia
es exactamente ese argumento, y es la parte que más puntos pierde.\visto{Examen dic. 2015 Ej. 2: mochila y principio de optimalidad}

\subsection{Qué espera la corrección}

Las soluciones oficiales de 2024 y 2025 traen la lista explícita de lo que se
evalúa. Es siempre la misma, sin importar el problema:\visto{2P 2024 Ej. 1 y 2P 2025 Ej. 3: criterios de corrección}

\begin{enumerate}[label=\textbf{\arabic*.}]
  \item \textbf{Definir la función.} Escribir $\OPT(\ldots)$ con su
  significado en palabras (``la ganancia máxima que se puede obtener desde el
  día $i$ hasta el día $n$'') y el \emph{rango} de cada parámetro. Si hay un
  parámetro de estado (la cuchilla instalada, los neumáticos usados), decir
  con precisión qué representa. Decir dónde queda la respuesta: $\OPT(n)$,
  $\OPT(1)$, $\min_X \OPT(n,X)$, $\max_j \OPT(j)$.
  \item \textbf{Casos base.} Los subproblemas que se resuelven sin consultar
  otros. Hacen falta tantos como ``saltos'' tenga la recurrencia: si
  $\OPT(i)$ usa $\OPT(i+7)$, hay que definir $\OPT(i)=0$ para
  $n<i\le n+7$.
  \item \textbf{Recurrencia justificada.} Se distingue por casos según la
  última (o la primera) decisión y se explica de dónde sale cada término. Lo
  ideal es mostrar las dos desigualdades: que cada término es el valor de
  alguna solución factible ($\ge$ o $\le$ según se maximice o minimice) y que
  el óptimo cae en alguno de los casos.
  \item \textbf{Algoritmo iterativo.} Una tabla con una celda por
  subproblema; se inicializan los casos base; se recorre en un orden tal que,
  al calcular una celda, las que usa ya estén calculadas; se devuelve la celda
  (o la combinación de celdas) que contiene la respuesta.
  \item \textbf{Reconstrucción.} Recorrer la tabla en sentido opuesto al de
  llenado; en cada paso decidir \emph{con la misma ecuación} con la que se
  calculó la celda, mantener el estado elegido y saltar las celdas que la
  decisión vuelve incompatibles.
  \item \textbf{Orden del tiempo de ejecución.} Cantidad de subproblemas por
  costo de calcular cada uno, más los precálculos (ordenar, calcular
  $p(j)$, errores $e_{i,j}$, etc.).
\end{enumerate}

\begin{resultado}[Plantilla de la justificación de una recurrencia]
Sea $S$ una solución óptima de la instancia $I$ (con valor $\OPT(I)$).
Mirando la última decisión de $S$ hay $k$ casos posibles. En el caso $c$, lo
que queda de $S$ al sacar esa decisión es una solución factible de la
subinstancia $I_c$, y tiene que ser óptima para ella: si hubiera una mejor,
le agrego la decisión y obtengo una solución de $I$ mejor que $S$. Entonces
$\OPT(I)=\text{costo}(c)+\OPT(I_c)$ para el caso en que cae $S$. Como
recíprocamente cada $\text{costo}(c)+\OPT(I_c)$ es el valor de una solución
factible de $I$, el óptimo es el mejor de los $k$ valores.
\end{resultado}

La versión rigurosa de este argumento, con las dos desigualdades por
separado, está en la solución oficial del ladrón de casas
(sección~\ref{pd:sec:wis}); vale la pena copiarle el formato.

\subsection{Memoización e iterativo}

Escribir la recurrencia como una función recursiva directa suele ser
exponencial, porque el mismo subproblema se resuelve muchas veces. Hay dos
formas de evitarlo, equivalentes en orden:
\begin{itemize}
  \item \textbf{Memoización (de arriba hacia abajo).} La función recursiva
  guarda cada resultado en un arreglo global $M$ y, antes de calcular,
  consulta si ya lo tiene.
  \item \textbf{Iterativa (de abajo hacia arriba).} Se llena la tabla en un
  orden compatible con las dependencias. Es lo que piden \emph{siempre} los
  parciales (``dé un algoritmo iterativo'').
\end{itemize}

El ejemplo canónico es Fibonacci. La versión dividir y conquistar
$F(n)=F(n-1)+F(n-2)$ hace $S(n)=1+S(n-1)+S(n-2)$ sumas, con $S(0)=S(1)=0$,
lo que da $S(n)=F(n)-1$, exponencial. La versión iterativa hace $n-1$
sumas y, si solo guarda los dos últimos valores, usa espacio
$\Oh(1)$.\visto{Examen feb. 2017 Preg. 3: Fibonacci con DyC y con PD}

\pdalgo{Fibonacci con PD en espacio constante}
\begin{algorithmic}[1]
\Function{FibPD}{$n$}
  \State $a \gets 0$;\ $b \gets 1$ \Comment{$a=F(i-1)$, $b=F(i)$}
  \For{$i=2$ \textbf{hasta} $n$}
    \State $c \gets a+b$;\ $a \gets b$;\ $b \gets c$
  \EndFor
  \State \Return $b$
\EndFunction
\end{algorithmic}

La moraleja que se pide explicar: la diferencia de crecimiento se debe a la
\emph{superposición de subproblemas}; la PD resuelve cada uno una sola vez.
El mismo truco de guardar solo la última fila (o las últimas $k$ celdas)
reduce la memoria de cualquier tabla cuando solo se pide el valor óptimo y
no la solución.\visto{2P 2017 Ej. 1: reducir la memoria de una tabla}

\begin{receta}[estructura de cualquier respuesta de PD]
\begin{enumerate}[leftmargin=1.2em]
  \item ``Para $\ldots$ (rango), definimos $\OPT(\ldots)$ como $\ldots$ (en
  palabras). La respuesta es $\ldots$''.
  \item ``Casos base: $\ldots$ porque $\ldots$''.
  \item ``Para $\ldots$, consideramos una solución óptima $S$ de $\ldots$.
  O bien $\ldots$ (caso 1), y entonces $\ldots$ es óptima para $\ldots$, con
  valor $\ldots$; o bien $\ldots$ (caso 2) $\ldots$. Por lo tanto
  $\OPT(\ldots)=\max\{\ldots\}$''.
  \item Algoritmo: casos base, ciclos en el orden correcto, una línea que
  aplica la recurrencia, \textbf{retornar}.
  \item Reconstrucción: \textbf{mientras} (no \textbf{para}) que recorre al
  revés y repite las comparaciones de la recurrencia.
  \item Orden: ``hay $\Theta(\ldots)$ subproblemas y cada uno se calcula en
  $\Oh(\ldots)$''.
\end{enumerate}
\end{receta}

\begin{trampa}
\begin{itemize}[leftmargin=1.2em]
  \item Una ``recurrencia'' que decide localmente lo mejor en cada paso (y
  por lo tanto podría construir la solución en la misma pasada) es un
  greedy disfrazado; casi siempre es incorrecta y se corrige como tal.
  \item Falta el parámetro que hace falta para que la decisión dependa de lo
  anterior (qué cuchilla está instalada, cuántas vueltas tienen los
  neumáticos).
  \item El significado de $\OPT$ cambia entre el lado izquierdo y el derecho
  de la ecuación.
  \item Casos base insuficientes: la recurrencia consulta $\OPT(i+2)$ o
  $\OPT(i-d)$ fuera del rango definido.
  \item La reconstrucción compara entre sí los valores de la tabla en cada
  posición (``si $\OPT(i,A)<\OPT(i,B)$ elijo $A$'') en lugar de seguir la
  decisión que produjo el valor.
  \item Algoritmo recursivo sin memoización cuando se pide iterativo.
  \item Orden ``$\Oh(nW)$, polinomial'': es \emph{pseudo}polinomial.
\end{itemize}
\end{trampa}

\subsection{Catálogo de patrones}

Casi todos los ejercicios de examen caen en alguno de estos moldes. Reconocer
el molde es la mitad del ejercicio: te dice qué parámetros usar, qué casos
base escribir y en qué orden llenar la tabla.

{\small
\begin{tabularx}{\textwidth}{@{}>{\raggedright\arraybackslash}p{2.9cm}>{\raggedright\arraybackslash}p{3.4cm}>{\raggedright\arraybackslash}X@{}}
\toprule
\textbf{Patrón} & \textbf{Subproblema} & \textbf{Ejemplos de examen}\\
\midrule
Prefijo con decisión binaria & $\OPT(i)$ sobre $1..i$; tomo $i$ o no & ladrón, baldosas, edificios, intervalos ponderados\\
Sufijo con saltos & $\OPT(i)$ sobre $i..n$; salto a $i+k$ & tarifa plana, pedidos, ravioles, dinosaurio\\
Prefijo con estado & $\OPT(i,X)$, $X$ = última elección & neumáticos A/B, cuchillas, F1\\
Segmentación & $\OPT(j)$; elijo dónde empieza el último tramo & fragmentos, prefijos, palíndromos\\
Último elemento & $\OPT(j)$ = mejor que termina en $j$ & subsecuencia creciente\\
Capacidad extra & $\OPT(i,w)$ con $w\le W$ & mochila, cambio, varillas\\
Dos secuencias & $\OPT(i,j)$ sobre prefijos de $x$ e $y$ & edición, alineación, ejecución paralela\\
Intervalo & $\OPT(i,j)$ sobre $x_i..x_j$ & matrices, paréntesis, palíndromo\\
Grilla & $\OPT(i,j)$ por casilla & sargento, minero, recolector\\
DAG o árbol & un valor por vértice, en orden topológico & camino más largo, conteo, conjunto independiente\\
Cantidad de aristas & $\OPT(i,v)$, a lo sumo $i$ aristas & Bellman--Ford, routers\\
\bottomrule
\end{tabularx}}

Las secciones que siguen recorren los moldes en ese orden. En cada una está la
teoría del libro, el pseudocódigo modelo y los ejercicios que se tomaron.
Las referencias (6.x) siguen la numeración del capítulo 6 de Kleinberg y
Tardos; si tu edición corre algún número, lo que importa es el enunciado.

% =====================================================================
\section{Intervalos ponderados y prefijos con decisión binaria}\label{pd:sec:wis}
% =====================================================================

\subsection{Weighted interval scheduling}

Hay $n$ pedidos; el pedido $j$ ocupa el intervalo $[s_j,f_j)$ y tiene valor
$v_j>0$. Se busca un subconjunto de pedidos disjuntos dos a dos de valor
máximo. Greedy por fin más temprano ya no sirve (solo sirve con valores
iguales). Se ordenan los pedidos por tiempo de fin,
$f_1\le f_2\le\dots\le f_n$, y se define $p(j)$ como el mayor índice $i<j$
tal que $i$ y $j$ son compatibles ($p(j)=0$ si no hay ninguno). Para
$0\le j\le n$, $\OPT(j)$ es el valor óptimo usando solo los pedidos
$1,\dots,j$; $\OPT(0)=0$.

Para el pedido $j$ hay dos casos: si está en la solución óptima $\mathcal{O}_j$,
ninguno de $p(j)+1,\dots,j-1$ puede estar (se solapan con $j$) y lo que queda
es una solución óptima para $1..p(j)$; si no está, $\mathcal{O}_j$ es una
solución óptima para $1..j-1$.

\begin{resultado}[(6.1)]
$\OPT(j)=\max\bigl(v_j+\OPT(p(j)),\ \OPT(j-1)\bigr)$.
\end{resultado}

\begin{resultado}[(6.2)]
El pedido $j$ pertenece a alguna solución óptima del conjunto
$\{1,\dots,j\}$ si y solo si $v_j+\OPT(p(j))\ge\OPT(j-1)$.
\end{resultado}

Implementada tal cual en forma recursiva (\textproc{Compute-Opt}), la
recurrencia es correcta (6.3, por inducción en $j$) pero exponencial: si
$p(j)=j-2$ para todo $j$, la cantidad de llamadas crece como Fibonacci. Con
memoización o en forma iterativa, cada valor se calcula una vez.

\pdalgo{Intervalos ponderados: memoización, iterativo y reconstrucción}
\begin{algorithmic}[1]
\Function{M-Compute-Opt}{$j$}
  \If{$j=0$} \Return $0$
  \ElsIf{$M[j]$ está definido} \Return $M[j]$
  \Else
    \State $M[j]\gets\max\bigl(v_j+\Call{M-Compute-Opt}{p(j)},\ \Call{M-Compute-Opt}{j-1}\bigr)$
    \State \Return $M[j]$
  \EndIf
\EndFunction
\Statex
\Function{Iterative-Compute-Opt}{\,}
  \State $M[0]\gets 0$
  \For{$j=1$ \textbf{hasta} $n$}
    \State $M[j]\gets\max\bigl(v_j+M[p(j)],\ M[j-1]\bigr)$
  \EndFor
\EndFunction
\Statex
\Function{Find-Solution}{$j$}
  \If{$j=0$} \Return $\emptyset$
  \ElsIf{$v_j+M[p(j)]\ge M[j-1]$}
    \State \Return $\{j\}\cup\Call{Find-Solution}{p(j)}$
  \Else\ \Return \Call{Find-Solution}{$j-1$}
  \EndIf
\EndFunction
\end{algorithmic}

\begin{resultado}[(6.4) y (6.5)]
Con los pedidos ordenados por fin y los $p(j)$ calculados,
\textproc{M-Compute-Opt}$(n)$ corre en $\Oh(n)$: cada llamada que no
devuelve un valor ya guardado llena una celda nueva y genera dos llamadas, y
hay $n$ celdas. Dado el arreglo $M$, \textproc{Find-Solution} devuelve una
solución óptima en $\Oh(n)$. Ordenar y calcular los $p(j)$ (con búsqueda
binaria) cuesta $\Oh(n\log n)$.
\end{resultado}

\subsection{El ladrón: la solución modelo}

Este es el ejercicio que conviene tener memorizado, porque la solución
oficial es la referencia de cómo se justifica una recurrencia. Las casas
$c_1,\dots,c_n$ están en fila, la casa $c_i$ tiene valor $v_i>0$ y no se
pueden robar dos contiguas. Es intervalos ponderados con $p(i)=i-2$: cada
casa se ``solapa'' solo con sus vecinas.\visto{2P 2019 Ej. 2: ladrón de casas contiguas}

\begin{problema}{Ladrón de casas no contiguas (2P 2019, Ej. 2)}
\textbf{Subproblema.} Para $0\le i\le n$, $\OPT(i)$ es la ganancia máxima
robando casas de $\{c_1,\dots,c_i\}$. La respuesta es $\OPT(n)$.

\textbf{Recurrencia.} $\OPT(0)=0$, $\OPT(1)=v_1$ y, para $i>1$,
\[\OPT(i)=\max\{\OPT(i-1),\ \OPT(i-2)+v_i\}.\]

\textbf{Justificación ($\ge$).} Sean $S_1\subseteq\{c_1..c_{i-1}\}$ y
$S_2\subseteq\{c_1..c_{i-2}\}$ compatibles de valores $\OPT(i-1)$ y
$\OPT(i-2)$. Tanto $S_1$ como $S_2\cup\{c_i\}$ son compatibles dentro de
$\{c_1..c_i\}$, así que $\OPT(i)$ es al menos el máximo de ambos valores.

\textbf{Justificación ($\le$).} Sea $S$ compatible con
$v(S)=\OPT(i)$. Si $c_i\notin S$, entonces $S\subseteq\{c_1..c_{i-1}\}$ y
$v(S)\le\OPT(i-1)$. Si $c_i\in S$, entonces $c_{i-1}\notin S$ y
$S\setminus\{c_i\}$ es compatible dentro de $\{c_1..c_{i-2}\}$, así que
$v(S)-v_i\le\OPT(i-2)$. En ambos casos $v(S)$ es a lo sumo el máximo.

\textbf{Algoritmo.} $M[0]\gets0$, $M[1]\gets v_1$; para $i=2..n$ aplicar la
recurrencia; devolver $M[n]$. Tiempo $\Oh(n)$.

\textbf{Reconstrucción.} $i\gets n$; \emph{mientras} $i\ge1$: si
$M[i]=M[i-1]$, $i\gets i-1$; si no, agregar $c_i$ e $i\gets i-2$.
\end{problema}

El mismo enunciado, con baldosas en lugar de casas (no se pueden retirar dos
baldosas seguidas), se tomó en diciembre de 2017. La solución guarda un
arreglo booleano de \emph{traceback} con la decisión tomada en cada $i$, que
es una alternativa válida a recalcular la comparación al
reconstruir.\visto{Examen dic. 2017 Prob. B: baldosas no contiguas}

Si la incompatibilidad alcanza a más vecinos, cambia solo el salto. En el
parcial de 2015 la restricción era una distancia mínima $d$ entre edificios,
así que el salto es $p(i)=i-d$; el ejercicio pedía también refutar un greedy
y corregir una implementación dada.\visto{2P 2015 Ej. 1: edificios a distancia mínima $d$}

\begin{problema}{Edificios separados por al menos $d$ cuadras (2P 2015, Ej. 1)}
La cuadra $i$ aloja $h_i$ habitantes (0 si no hay solicitud) y dos edificios
aceptados deben estar a distancia $\ge d$. Con $T(i)$ el máximo para las
cuadras $1..i$:
\[
T(0)=0,\qquad T(i)=\max\bigl\{h_i+T(\max\{0,i-d\}),\ T(i-1)\bigr\},\quad i\ge1.
\]
Es intervalos ponderados con $p(i)=\max\{0,i-d\}$. Se pedía además mostrar
que el greedy ``aceptar en orden decreciente de habitantes si es compatible''
falla: con $d=2$ y $h=(2,3,2)$ elige $\{2\}$ (valor 3) y el óptimo es
$\{1,3\}$ (valor 4).
\end{problema}

\subsection{Sufijos con saltos}

Muchas veces es más natural mirar la \emph{primera} decisión: $\OPT(i)$ es el
óptimo para los días $i,i+1,\dots,n$, se llena con $i$ decreciente y la
respuesta es $\OPT(1)$. Si la decisión del día $i$ ``consume'' $k$ días, la
recurrencia salta a $\OPT(i+k)$ y hacen falta casos base hasta
$n+k$.\visto{2P 2021 Ej. 1: tarifa plana de 7 días}

\begin{problema}{Tarifa por consumo o tarifa plana (2P 2021, Ej. 1)}
El día $i$ se paga $t_iK_c$ por consumo, o se contrata un bloque de 7 días
consecutivos por $K_p$ (aunque queden menos de 7 días). Para
$1\le i\le n+7$, $\OPT(i)$ es el costo óptimo de los días $i..n$:
\begin{gather*}
\OPT(i)=0\qquad (n<i\le n+7),\\
\OPT(i)=\min\{t_iK_c+\OPT(i+1),\ K_p+\OPT(i+7)\}\qquad (1\le i\le n).
\end{gather*}
El día $i$ solo puede pagarse de esas dos formas; con la plana quedan pagos
hasta $\min\{i+6,n\}$. Se llena con $i=n$ bajando hasta 1. Reconstrucción:
$i\gets1$; mientras $i\le n$, si $\OPT(i)=K_p+\OPT(i+7)$ agregar $i$ a $P$ e
$i\gets i+7$, si no $i\gets i+1$.
\end{problema}

\begin{problema}{Ravioles y sorrentinos (examen dic. 2024, Ej. 2)}
Cada día se fabrica un solo tipo de pasta; si el día $i$ se hacen
sorrentinos, los días $i+1$ e $i+2$ solo pueden ser ravioles. Las ventas
estimadas son $r_i$ y $s_i$ unidades, con ganancias unitarias $g_r$ y $g_s$.
Con $\OPT(i)$ la ganancia máxima de los días $i..n$ y $r_{n+1}=r_{n+2}=0$:
\begin{gather*}
\OPT(n+d)=0\qquad (1\le d\le3),\\
\OPT(i)=\max\bigl\{s_ig_s+(r_{i+1}+r_{i+2})g_r+\OPT(i+3),\ r_ig_r+\OPT(i+1)\bigr\}.
\end{gather*}
Si se elige sorrentinos el día $i$, los dos días siguientes quedan fijados y
lo que sigue desde $i+3$ no tiene restricciones. Reconstrucción con
\textbf{mientras}: si $\OPT(i)=r_ig_r+\OPT(i+1)$ poner $R$ y avanzar 1; si no,
poner $S,R,R$ y avanzar 3.
\end{problema}

El de pedidos de soporte de 2025 es el mismo molde, pero con tres casos que
dependen de los datos y no de la decisión, y la solución oficial lista
errores típicos de la reconstrucción que conviene leer.\visto{2P 2025 Ej. 3: pedidos que duran uno o dos días}

\begin{problema}{Pedidos de soporte de uno o dos días (2P 2025, Ej. 3)}
El pedido del día $i$ requiere $r_i$ horas y da ganancia $g_i$; el día $i$
hay $d_i$ horas. Si $r_i\le d_i$ se hace en el día; si
$d_i<r_i\le d_i+d_{i+1}$ ocupa también el día $i+1$ (y se pierde el pedido
$i+1$); si $r_i>d_i+d_{i+1}$ se rechaza. El día $n$ solo se acepta si
$r_n\le d_n$. Con $\OPT(i)$ la ganancia máxima de los días $i..n$:
\[
\OPT(n+1)=0,\qquad \OPT(n)=\begin{cases}g_n & r_n\le d_n\\ 0 & r_n>d_n\end{cases}
\]
\[
\OPT(i)=\begin{cases}
\OPT(i+1) & r_i>d_i+d_{i+1}\\
g_i+\OPT(i+1) & r_i\le d_i\\
\max\{\OPT(i+1),\ g_i+\OPT(i+2)\} & d_i<r_i\le d_i+d_{i+1}
\end{cases}\qquad 1\le i<n.
\]
Se llena desde $n-1$ hasta 1. En el ejemplo de la letra la respuesta es
$G=11$.

\textbf{Reconstrucción.} Con \textbf{mientras}: si $\OPT(i)=\OPT(i+1)$, se
rechaza y se avanza 1; si no, se acepta, y se avanza 1 si $r_i\le d_i$ y 2 en
otro caso. Hay que decidir mirando también $r$ y $d$, no solo los valores de
la tabla.
\end{problema}

\begin{trampa}[reconstrucción en sufijos con saltos]
\begin{itemize}[leftmargin=1.2em]
  \item Usar \textbf{para} $i=1..n$: no permite saltar el día que quedó
  ocupado. Con $n=2$, $g=(20,10)$, $d=(1,1)$, $r=(2,1)$ la tabla da
  $\OPT(2)=10$, $\OPT(1)=20$ y un \textbf{para} acepta los dos pedidos.
  \item Comparar solo valores: puede valer $\OPT(i)=g_i+\OPT(i+2)$ aunque
  $r_i>d_i+d_{i+1}$ (pedido imposible). Con $n=3$, $g=(10,10,5)$,
  $d=(1,1,1)$, $r=(3,2,2)$ se tiene $\OPT(3)=0$, $\OPT(2)=10$, $\OPT(1)=10$,
  y el pedido 1 se aceptaría por error.
  \item Olvidar el caso base extra ($n+1$, o $n+1..n+7$) cuando la recurrencia
  salta más de un lugar.
\end{itemize}
\end{trampa}

\subsection{Recurrencias booleanas}

Si la pregunta es ``¿se puede?'', $\OPT$ es booleana y el $\max$ se vuelve
$\lor$. La reconstrucción es más simple que en optimización: basta ir a
cualquier sucesor que valga \texttt{true}.\visto{Examen feb. 2020 Ej. 3: dinosaurio que salta obstáculos}

\begin{problema}{Dinosaurio que avanza 1 o salta 4 (examen feb. 2020, Ej. 3)}
Hay $n$ baldosas, libres o con obstáculo ($B[i]$). Desde $i$ se va a $i+1$ o a
$i+4$; se gana al pasar más allá de $n$. Con $M(i)$ verdadero si se puede
ganar partiendo de $i$:
\begin{gather*}
M(i)=\texttt{true}\qquad (n<i\le n+4),\\
M(i)=\bigl(B[i]=\text{libre}\bigr)\wedge\bigl(M(i+1)\vee M(i+4)\bigr)\qquad (1\le i\le n).
\end{gather*}
Se calcula con $i$ decreciente. Si $M(1)$ es falso no hay solución; si no,
desde $i=1$, mientras $i\le n$: si $M(i+4)$ saltar, si no avanzar.
\end{problema}

\begin{receta}[prefijos y sufijos]
\begin{itemize}[leftmargin=1.2em]
  \item Si la letra sugiere ``desde el día $i$ hasta el $n$'', usá sufijos y
  llená con $i$ decreciente; si sugiere ``los primeros $i$'', prefijos.
  \item Escribí los casos base para todos los índices fuera de rango que
  aparecen en la recurrencia (o agregá ceros ficticios, $r_{n+1}=0$, y decilo).
  \item En cada término explicá en una línea qué decisión representa y por
  qué lo que queda es independiente de esa decisión.
  \item Si se pide comparar con un greedy, el contraejemplo de tres elementos
  casi siempre alcanza.
\end{itemize}
\end{receta}

% =====================================================================
\section{Prefijos con un parámetro de estado}
% =====================================================================

Cuando el costo de una decisión depende de la decisión anterior (cambiar de
cuchilla cuesta $T$, cambiar neumáticos cuesta $P$), $\OPT(i)$ no alcanza: al
escribir la recurrencia no sabés qué se eligió en $i-1$. Se agrega un
parámetro que guarda exactamente esa información, y la tabla pasa a tener
una fila por cada valor del estado. Este molde salió en el segundo parcial
de 2024 con la mitad del puntaje.\visto{2P 2024 Ej. 1: cuchillas A y B}

\begin{problema}{Cuchillas A y B (2P 2024, Ej. 1)}
Hay $n$ cortes en orden sobre materiales $m_1..m_n$; cortar $m$ con cuchilla
$X\in\{A,B\}$ cuesta $C_X(m)$ y cambiar de cuchilla cuesta $T$. Al empezar
está instalada la $A$. Escribimos $\pdbar{X}$ para la cuchilla opuesta.

\textbf{Subproblema.} Para $1\le i\le n$ y $X\in\{A,B\}$, $\OPT(i,X)$ es el
costo mínimo de hacer las primeras $i$ tareas \emph{terminando con la
cuchilla $X$} (la tarea $i$ se hace con $X$).

\textbf{Recurrencia.}
\[
\OPT(1,A)=C_A(m_1),\qquad \OPT(1,B)=C_B(m_1)+T,
\]
\[
\OPT(i,X)=C_X(m_i)+\min\bigl\{\OPT(i-1,X),\ \OPT(i-1,\pdbar{X})+T\bigr\},\quad 1<i\le n.
\]
Si la tarea $i$ se hace con $X$, se paga el desgaste $C_X(m_i)$; las
primeras $i-1$ tareas deben resolverse de forma óptima terminando con $X$
(sin cambio) o con $\pdbar{X}$ (y se paga el cambio). La respuesta es
$\min\{\OPT(n,A),\OPT(n,B)\}$. Tiempo $\Oh(n)$.

\textbf{Reconstrucción.} $X\gets\pdargmin_X\OPT(n,X)$; para $i=n$ bajando
hasta 2: $\mathit{SOL}[i]\gets X$; si $\OPT(i-1,\pdbar{X})+T<\OPT(i-1,X)$ entonces
$X\gets\pdbar{X}$. Al final $\mathit{SOL}[1]\gets X$.
\end{problema}

La solución oficial de este ejercicio es la más útil para estudiar, porque
analiza tres recurrencias incorrectas que aparecieron en las pruebas. Las
tres fallan por lo mismo: la información de qué cuchilla se usó en la tarea
anterior no está en el subproblema, o está con un significado distinto del
que se usa en el lado derecho.

\begin{trampa}[las recurrencias incorrectas del 2P 2024]
\begin{itemize}[leftmargin=1.2em]
  \item $\OPT(i)=\OPT(i-1)+\min\{C_A(m_i),T+C_B(m_i)\}$ ``si la tarea $i-1$
  se hizo con $A$'': no hay forma de saber eso desde $\OPT(i-1)$; es greedy.
  \item ``El último cambio fue en la tarea $k$'': para sumar $T$ hay que
  saber que la tarea $k-1$ se hizo con la otra cuchilla, y eso no está en
  $\OPT(k-1)$.
  \item $\OPT(i,X)$ con $X$ = ``cuchilla instalada al empezar la tarea $i$''
  pero usando $\OPT(i-1,X)$ como si $X$ fuera la cuchilla con la que se hizo
  $i-1$: el parámetro cambia de significado entre los dos lados.
  \item Reconstruir eligiendo en cada $i$ la cuchilla con menor $\OPT(i,X)$:
  con $T=5$, $C_A=(10,10)$, $C_B=(1,100)$ la tabla es
  $\OPT(\cdot,A)=(10,20)$, $\OPT(\cdot,B)=(6,106)$ y ese método devuelve
  $[B,A]$ de costo 21, cuando el óptimo es $[A,A]$ de costo 20.
\end{itemize}
\end{trampa}

El de neumáticos de diciembre de 2019 es idéntico salvo por el caso base: el
tipo de neumático de la largada se elige libremente, así que
$\OPT(1,X)=t(1,X)$ para ambos $X$, y la parada cuesta
$P$.\visto{Examen dic. 2019 Ej. 2: neumáticos A o B por vuelta}

\pdalgo{Neumáticos A/B: costo mínimo y solución}
\begin{algorithmic}[1]
\State $\OPT[1,X]\gets t(1,X)$ para $X\in\{A,B\}$
\For{$i=2$ \textbf{hasta} $n$}
  \ForAll{$X\in\{A,B\}$}
    \State $\OPT[i,X]\gets t(i,X)+\min\{\OPT[i-1,X],\ \OPT[i-1,\pdbar{X}]+P\}$
  \EndFor
\EndFor
\State $X\gets\pdargmin_{X}\OPT[n,X]$
\For{$i=n$ \textbf{bajando hasta} $2$}
  \State $\mathit{SOL}[i]\gets X$
  \If{$\OPT[i-1,\pdbar{X}]+P<\OPT[i-1,X]$} $X\gets\pdbar{X}$ \EndIf
\EndFor
\State $\mathit{SOL}[1]\gets X$;\ \Return $\min_X\OPT[n,X]$ y $SOL$
\end{algorithmic}

Un estado puede ser también un número. En la estrategia de Fórmula 1 de 2022
el tiempo de cada vuelta depende de cuántas vueltas tienen los neumáticos,
así que ese contador entra en el subproblema. La solución oficial lo vincula
con el ejercicio 6.9 del libro (procesamiento con reinicios).\visto{2P 2022 Ej. 1: paradas de Fórmula 1}

\begin{problema}{Paradas en boxes de Fórmula 1 (2P 2022, Ej. 1)}
$s_{i,j}$ es el tiempo de la vuelta $i$ con neumáticos de $j$ vueltas de uso
($1\le j\le i\le n$); parar al final de una vuelta cuesta $P$. Con
$\OPT(i,j)$ el tiempo mínimo desde la vuelta $i$ hasta la $n$ si los
neumáticos llevan $j$ vueltas al correr la vuelta $i$:
\begin{gather*}
\OPT(n,j)=s_{n,j},\\
\OPT(i,j)=s_{i,j}+\min\bigl\{P+\OPT(i+1,1),\ \OPT(i+1,j+1)\bigr\}\qquad(1\le j\le i<n).
\end{gather*}
Respuesta $\OPT(1,1)$; se llena con $i$ decreciente y, para cada $i$, todos
los $j\le i$: $\Oh(n^2)$. Reconstrucción: $j\gets1$; para $i=1..n-1$, si
$\OPT(i,j)=s_{i,j}+P+\OPT(i+1,1)$ se para en $i$ y $j\gets1$, si no
$j\gets j+1$.

\textbf{Alternativa.} $\OPT'(i,j)$ = tiempo óptimo de las vueltas $1..i$
terminando con neumáticos de $j$ vueltas: $\OPT'(i,j)=\OPT'(i-1,j-1)+s_{i,j}$
si $j>1$, y $\OPT'(i,1)=P+s_{i,1}+\min_k\OPT'(i-1,k)$; respuesta
$\min_k\OPT'(n,k)$.
\end{problema}

En el programa anterior a 2018 apareció la versión ``dos agentes'': Ana y
Luis marcan los lugares $l_1,\dots,l_n$ en orden y en cada paso uno de los
dos va al siguiente lugar. El estado es dónde quedó el que no se movió.\visto{2P 2010 Ej. 2: dos corredores marcan lugares}

\begin{problema}{Dos integrantes marcan lugares (2P 2010, Ej. 2)}
$f(l,k)$, con $k<l$, es la distancia que falta recorrer si ya se marcaron
los lugares hasta $l_l$, uno de los integrantes está en $l_l$ y el otro en
$l_k$. Caso base: $f(n,k)=0$. En general, el lugar $l_{l+1}$ lo marca el que
está en $l_l$ o el que está en $l_k$:
\[
f(l,k)=\min\bigl\{d(l_l,l_{l+1})+f(l+1,k),\ d(l_k,l_{l+1})+f(l+1,l)\bigr\}.
\]
Hay $\Oh(n^2)$ subproblemas de costo $\Oh(1)$.
\end{problema}

\begin{receta}[cuándo agregar un estado]
Preguntate: ``para decidir lo que hago en el paso $i$, ¿qué necesito saber
de lo que pasó antes?''. Eso, y nada más, es el parámetro extra. Escribí
explícitamente si $X$ es la elección hecha \emph{en} el paso $i$ o la que
estaba vigente \emph{al empezar} el paso $i$, y usá el mismo significado en
ambos lados de la ecuación. En la reconstrucción llevá el estado en una
variable y actualizalo con la misma comparación de la recurrencia.
\end{receta}

% =====================================================================
\section{Segmentación: elegir dónde empieza el último tramo}
% =====================================================================

\subsection{Mínimos cuadrados segmentados}

Hay puntos $p_1,\dots,p_n$ ordenados por abscisa. Se quiere partirlos en
tramos consecutivos y ajustar una recta a cada tramo; el costo es la suma de
los errores cuadráticos $e_{i,j}$ de cada tramo $p_i..p_j$ más una constante
$C$ por tramo. La decisión ya no es binaria: el último tramo empieza en
algún $i$ entre 1 y $n$, y hay que probar todos (decisión \emph{de múltiples
opciones}).

\begin{resultado}[(6.6)]
Si el último tramo de la partición óptima es $p_i,\dots,p_n$, entonces
$\OPT(n)=e_{i,n}+C+\OPT(i-1)$.
\end{resultado}

\begin{resultado}[(6.7)]
Para los puntos $p_1,\dots,p_j$,
$\OPT(j)=\min_{1\le i\le j}\bigl(e_{i,j}+C+\OPT(i-1)\bigr)$, con
$\OPT(0)=0$; el tramo $p_i..p_j$ se usa en una solución óptima del
subproblema si y solo si el mínimo se alcanza en $i$.
\end{resultado}

\pdalgo{Segmented-Least-Squares y reconstrucción}
\begin{algorithmic}[1]
\State calcular $e_{i,j}$ para todo $1\le i\le j\le n$
\State $M[0]\gets 0$
\For{$j=1$ \textbf{hasta} $n$}
  \State $M[j]\gets\min_{1\le i\le j}\bigl(e_{i,j}+C+M[i-1]\bigr)$
\EndFor
\Statex
\Function{Find-Segments}{$j$}
  \If{$j>0$}
    \State hallar $i$ que minimiza $e_{i,j}+C+M[i-1]$
    \State emitir el tramo $p_i..p_j$;\ \Call{Find-Segments}{$i-1$}
  \EndIf
\EndFunction
\end{algorithmic}

Con los errores precalculados, el ciclo principal es $\Oh(n^2)$ ($n$
subproblemas, $\Oh(n)$ cada uno). Calcular todos los $e_{i,j}$ de forma
directa cuesta $\Oh(n^3)$ y puede bajarse a $\Oh(n^2)$ con sumas
acumuladas. En los parciales el costo del tramo suele darse como función
calculable en $\Oh(1)$.\visto{2P 2020 Ej. 2: cortar una secuencia en fragmentos}

\begin{problema}{Fragmentos con costo máximo menos mínimo (2P 2020, Ej. 2)}
Partir $x_1..x_n$ en fragmentos consecutivos; cada fragmento cuesta
$d_{i,j}=\max-\min$ de sus elementos y cada \emph{corte} cuesta $c$ (no cada
fragmento). Con $\OPT(j)$ el costo mínimo del prefijo $x_1..x_j$:
\[
\OPT(j)=\min\Bigl\{d_{1,j},\ c+\min_{2\le i\le j}\bigl(\OPT(i-1)+d_{i,j}\bigr)\Bigr\},\qquad 1\le j\le n,
\]
con la convención $\min\emptyset=\infty$ (para $j=1$ queda $d_{1,1}=0$). El
primer término es ``un solo fragmento''; el segundo, ``el último fragmento
empieza en $i>1$'', que agrega un corte. Orden $\Oh(n^2)$.

\textbf{Posiciones de inicio.} $j\gets n$; mientras $j>0$: si
$\OPT(j)=d_{1,j}$ entonces $i\gets1$, si no buscar $i\ge2$ con
$\OPT(j)=c+\OPT(i-1)+d_{i,j}$; marcar $p[i]\gets\texttt{true}$;
$j\gets i-1$.
\end{problema}

Ojo con la diferencia respecto del libro: allá la constante se paga por
tramo, acá por corte. Una forma limpia de unificar es definir
$\OPT(0)=-c$ y usar $\OPT(j)=\min_{1\le i\le j}(c+\OPT(i-1)+d_{i,j})$. En
febrero de 2024 se tomó otra variante, en la que el ``tramo'' es un prefijo
de una cadena fija.\visto{Examen feb. 2024 Ej. 2: concatenación de prefijos}

\begin{problema}{Concatenación de prefijos (examen feb. 2024, Ej. 2)}
Se quiere construir $Y=y_1..y_n$ como concatenación de prefijos no nulos de
$S=s_1..s_n$ que coincida con $X$ en la mayor cantidad de posiciones. Si el
último prefijo cubre $x_{i+1}..x_j$, aporta
$c(i,j)=|\{k: i<k\le j,\ x_k=s_{k-i}\}|$ coincidencias:
\[
\OPT(0)=0,\qquad \OPT(j)=\max_{0\le i<j}\bigl(\OPT(i)+c(i,j)\bigr),\quad 1\le j\le n.
\]
Precalcular $c$ ingenuamente cuesta $\Oh(n^3)$; con
$c(i,j)=c(i,j-1)+[x_j=s_{j-i}]$ baja a $\Oh(n^2)$, igual que el ciclo
principal.
\end{problema}

La segmentación en palíndromos del examen de febrero de 2026 es el mismo
molde con la condición de que cada tramo sea palíndromo. La solución oficial
la presenta explícitamente como ``parecida a mínimos cuadrados
segmentados''.\visto{Examen feb. 2026 Ej. 3: segmentación en palíndromos}

\begin{problema}{Segmentación mínima en palíndromos (examen feb. 2026, Ej. 3)}
Con $\mathit{esP}(i,j)$ verdadero si $a_i..a_j$ es palíndromo y $\OPT(j)$ el tamaño
mínimo de una segmentación en palíndromos de $a_1..a_j$:
\[
\OPT(0)=0,\qquad \OPT(j)=1+\min\{\OPT(i-1):\ 1\le i\le j,\ \mathit{esP}(i,j)\}.
\]
El rango de $i$ tiene que incluir $i=1$ (toda la cadena es un palíndromo) e
$i=j$ (tramos de largo 1, que siempre existen, así que el conjunto nunca es
vacío). Reconstrucción: $j\gets n$; mientras $j>0$, elegir $i$ con
$\mathit{esP}(i,j)$ y $\OPT(j)=1+\OPT(i-1)$, agregar $i$ al inicio de $S$,
$j\gets i-1$. Orden $\Oh(n^2)$ con $\mathit{esP}$ precalculada (que es a su vez una
PD sobre intervalos, sección~\ref{pd:sec:intervalos}).

El greedy ``el palíndromo inicial más largo posible'' falla: para
\texttt{dqdqdmmmmdq} da $\{\texttt{dqdqd},\texttt{mmmm},\texttt{d},\texttt{q}\}$
y el óptimo es $\{\texttt{dqd},\texttt{qdmmmmdq}\}$. Tomar el palíndromo
final más largo también falla.
\end{problema}

\begin{trampa}[segmentación]
\begin{itemize}[leftmargin=1.2em]
  \item Excluir $i=1$ o $i=j$ del rango del mínimo, dejando afuera ``un solo
  tramo'' o ``tramos de largo 1''.
  \item Sumar la constante también al primer tramo cuando se cobra por corte.
  \item Escribir $\OPT(j)=\OPT(j-1)+\ldots$: eso decide solo si $j$ empieza o
  no un tramo, sin saber dónde empezó el anterior.
\end{itemize}
\end{trampa}

% =====================================================================
\section{Subsecuencia creciente más larga}
% =====================================================================

Otra decisión de múltiples opciones: para que el subproblema tenga
estructura, se fija el \emph{último} elemento. Para $1\le j\le n$,
$\OPT(j)$ es el largo máximo de una subsecuencia creciente de $x_1..x_j$ que
\emph{termina en $x_j$}. Si la subsecuencia tiene largo mayor a 1, su
penúltimo elemento es algún $x_i$ con $i<j$ y $x_i<x_j$, y lo anterior es una
subsecuencia creciente óptima que termina en $x_i$.\visto{Examen dic. 2023 Ej. 2: subsecuencia creciente más larga}

\begin{problema}{Subsecuencia creciente más larga (dic. 2023; 2P 2016; feb. 2018)}
Sea $S_j=\{i: 1\le i<j,\ x_i<x_j\}$. Entonces
\[
\OPT(j)=\begin{cases}1 & S_j=\emptyset\\ 1+\max\{\OPT(i): i\in S_j\} & S_j\ne\emptyset\end{cases}
\]
y la respuesta es $\max_{1\le j\le n}\OPT(j)$, no $\OPT(n)$. Alternativa de
2018: $G_0=0$, $G_i=\max\{G_{i-1},L_i\}$ con $L_i=\OPT(i)$, y la respuesta es
$G_n$. Orden $\Oh(n^2)$.

\textbf{Reconstrucción.} $j\gets\pdargmax_j\OPT(j)$; $P\gets(x_j)$; para
$i=j-1$ bajando hasta 1: si $x_i<x_j$ y $\OPT(i)+1=\OPT(j)$, agregar $x_i$ al
principio de $P$ y $j\gets i$. También puede guardarse un arreglo $T[j]$ con el
índice anterior elegido.
\end{problema}

En 2016 se pedía, además, explicar qué subproblemas hacen falta para resolver
$P_i$ (los $P_j$ con $j<i$ y $A_j<A_i$) y por qué la versión recursiva
directa es exponencial: en el peor caso, la secuencia creciente, la llamada
para $j$ se repite $2^{n-j}$ veces.\visto{2P 2016 Ej. 2: largos $L_i$ de subsecuencias crecientes}

En febrero de 2018 se agregó una reducción: apilar la mayor cantidad de
bandejas rectangulares, donde cada una debe ser más chica en ancho y largo
que la de abajo. Se ordenan por ancho y se busca la subsecuencia creciente más
larga de los largos.\visto{Examen feb. 2018 Prob. B: LIS y pila de bandejas}

\begin{trampa}[subsecuencia creciente]
Devolver $\OPT(n)$ (la mejor que termina en el último elemento) en lugar del
máximo sobre todos los $j$; definir $\OPT(j)$ como ``la más larga en
$x_1..x_j$'' y después usar que termina en $x_j$; olvidar que con elementos
repetidos la condición es estricta ($x_i<x_j$).
\end{trampa}

% =====================================================================
\section{Mochila y suma de subconjuntos: una capacidad extra}
% =====================================================================

\subsection{Suma de subconjuntos}

Hay $n$ ítems con pesos enteros $w_i$ y una cota $W$; se busca un
subconjunto $S$ con $\sum_{i\in S}w_i\le W$ lo más grande posible. Con
$\OPT(i)$ sola no alcanza: si el ítem $n$ entra en la solución, lo que
queda es un problema sobre $1..n-1$ \emph{con capacidad $W-w_n$}, que no es
ninguno de los $\OPT(i)$. Hay que agregar la capacidad restante como
parámetro: $\OPT(i,w)$ es el óptimo con los ítems $1..i$ y cota $w$, para
$0\le i\le n$ y $0\le w\le W$.

\begin{resultado}[(6.8)]
Si $w<w_i$ entonces $\OPT(i,w)=\OPT(i-1,w)$. Si no,
$\OPT(i,w)=\max\bigl(\OPT(i-1,w),\ w_i+\OPT(i-1,w-w_i)\bigr)$.
Casos base: $\OPT(0,w)=0$ para todo $w$.
\end{resultado}

\pdalgo{Subset-Sum$(n,W)$ y reconstrucción}
\begin{algorithmic}[1]
\State $M[0,w]\gets 0$ para $0\le w\le W$
\For{$i=1$ \textbf{hasta} $n$}
  \For{$w=0$ \textbf{hasta} $W$}
    \If{$w<w_i$} $M[i,w]\gets M[i-1,w]$
    \Else\ $M[i,w]\gets\max\bigl(M[i-1,w],\ w_i+M[i-1,w-w_i]\bigr)$
    \EndIf
  \EndFor
\EndFor
\State $S\gets\emptyset$;\ $w\gets W$
\For{$i=n$ \textbf{bajando hasta} $1$}
  \If{$w\ge w_i$ \textbf{y} $M[i,w]=w_i+M[i-1,w-w_i]$}
    \State $S\gets S\cup\{i\}$;\ $w\gets w-w_i$
  \EndIf
\EndFor
\State \Return $M[n,W]$ y $S$
\end{algorithmic}

\begin{resultado}[(6.9), (6.10) y (6.11)]
\textproc{Subset-Sum}$(n,W)$ calcula correctamente el óptimo en tiempo
$\Oh(nW)$. Dada la tabla, el conjunto óptimo se reconstruye en $\Oh(n)$. La
mochila (cada ítem tiene además valor $v_i$, y se maximiza
$\sum_{i\in S}v_i$ sujeto a $\sum_{i\in S}w_i\le W$) se resuelve igual en
$\Oh(nW)$ con
$\OPT(i,w)=\max\bigl(\OPT(i-1,w),\ v_i+\OPT(i-1,w-w_i)\bigr)$ si $w_i\le w$.
\end{resultado}

El orden $\Oh(nW)$ es \emph{pseudopolinomial}: es polinomial en el
\emph{valor} de $W$, pero $W$ se escribe con $\log W$ bits, así que puede ser
exponencial en el tamaño de la entrada. No contradice que suma de
subconjuntos sea $\NP$-completo (capítulo 8), y en un parcial hay que decirlo
así y no ``polinomial''.

\subsection{Mochila 0/1}

La mochila es el problema de PD más preguntado del programa viejo (hasta
2017), siempre con la misma estructura: plantear la recurrencia, justificar
por qué greedy no sirve, implementar y, a veces, comparar con la versión
fraccionaria.\visto{Examen jul. 2016 Prob. C: mochila de Martín}

\begin{problema}{Mochila 0/1 (jul. 2016 Prob. C; dic. 2010 Prob. 1; 2P 2008 Ej. 1)}
Productos con precio $d_k$ y peso $w_k$, un solo ejemplar de cada uno,
capacidad $W$. Con $g(k,w)$ la ganancia óptima usando los productos
$1..k$ con peso máximo $w$:
\[
g(0,w)=0,\qquad
g(k,w)=\begin{cases}g(k-1,w) & w<w_k\\ \max\{g(k-1,w),\ g(k-1,w-w_k)+d_k\} & w\ge w_k.\end{cases}
\]
Se llena por filas $k=1..n$ y, dentro de cada fila, $w=0..W$; respuesta
$g(n,W)$.

\textbf{Greedy no sirve.} Con $W=10$ y ítems $(w,d)=(6,7),(5,5),(5,5)$,
elegir por mayor $d/w$ toma el primero y obtiene 7; el óptimo es 10.

\textbf{Versión fraccionaria.} Si se pueden llevar fracciones, greedy por
mayor $d_k/w_k$ es óptimo, en $\Oh(n\log n)$.
\end{problema}

La misma estructura aparece disfrazada: espacios de publicidad con precio y
rating y un presupuesto (2012, donde también se pedía mostrar que el greedy
por cociente falla) y la definición formal del problema más el principio de
optimalidad (diciembre de 2015).\visto{Examen feb. 2012 Ej. 2: publicidad con presupuesto}

\subsection{Ítems ilimitados: cambio de monedas}

Si cada ítem se puede usar cualquier cantidad de veces, al usar el ítem $i$
\emph{no} se pasa a $i-1$: se queda en la misma fila con menos capacidad.
Ese es el único cambio, y es el error más común.\visto{Examen dic. 2020 Ej. 3: cambio con menos billetes}

\begin{problema}{Cambio con la menor cantidad de billetes (dic. 2020 Ej. 3; jul. 2020 Ej. 3)}
Valores $1=w_1<w_2<\dots<w_n$, cantidad ilimitada, monto $C$. Con
$\OPT(i,R)$ la menor cantidad de billetes con valores en $\{w_1..w_i\}$ que
suman $R$:
\[
\OPT(i,0)=0,\qquad \OPT(1,R)=R,
\]
\[
\OPT(i,R)=\begin{cases}\OPT(i-1,R) & R<w_i\\ \min\{\OPT(i-1,R),\ 1+\OPT(i,R-w_i)\} & R\ge w_i.\end{cases}
\]
O no se usa ningún billete $w_i$, o se usa al menos uno y el resto
$R-w_i$ se paga de forma óptima con los mismos valores. Orden $\Oh(nC)$.

\textbf{Greedy falla.} Con $B=\{1,3,4\}$ y $C=6$, tomar siempre el mayor
billete da $4+1+1$ (3 billetes) y el óptimo es $3+3$. En el problema de las
varillas (longitudes $1,4,6$) las respuestas para $4,8,14,20$ son $1,2,3,4$
varillas; por ejemplo $14=6+4+4$ y $20=6+6+4+4$.

\textbf{Reconstrucción.} $i\gets n$, $R\gets C$; mientras $R>0$: si $i=1$,
usar $R$ billetes de 1; si no, si $R<w_i$ o $\OPT(i-1,R)<1+\OPT(i,R-w_i)$,
$i\gets i-1$; si no, usar un billete $w_i$ y $R\gets R-w_i$.
\end{problema}

También se puede usar una sola dimensión,
$\OPT(R)=1+\min\{\OPT(R-w_i): w_i\le R\}$ con $\OPT(0)=0$, que es
correcta porque el orden de los billetes no importa. La ventaja de la versión
con dos índices es que es la que sugiere la letra y la reconstrucción es más
directa.\visto{Examen jul. 2020 Ej. 3: varillas con mínimas soldaduras}

El greedy de monedas (tomar la de mayor valor que entra, tantas veces como
se pueda) es óptimo para sistemas como el uruguayo, pero no en general; en
febrero de 2018 se pedía identificar la técnica y mostrar que puede no
encontrar la solución óptima. La PD de arriba es la que resuelve el caso
general.\visto{Examen feb. 2018 Preg. 2: greedy de monedas}

\subsection{Variantes con cantidades acotadas o dos restricciones}

Si de cada ítem hay un stock limitado, la decisión pasa a ser \emph{cuántas}
unidades del último ítem usar, y se toma el mínimo (o el máximo) sobre esa
cantidad. Si además la capacidad debe llenarse exactamente, las celdas
imposibles valen $+\infty$ y la respuesta puede no existir.\visto{2P 2012 Ej. 1: carga exacta con stock}

\begin{problema}{Carga exacta con stock acotado (2P 2012, Ej. 1)}
Productos $0..n$ de volumen $T[p]$ con stock $Q[p]$; llenar \emph{exactamente}
la capacidad $C$ con la menor cantidad de unidades. Con $f(p,c)$ el mínimo
usando los productos $0..p$ para volumen exacto $c$ ($+\infty$ si es
imposible), una recurrencia posible es
\begin{gather*}
f(-1,0)=0,\qquad f(-1,c)=+\infty\quad (c>0),\\
f(p,c)=\min_{0\le q\le\min\{Q[p],\lfloor c/T[p]\rfloor\}}\bigl(q+f(p-1,c-q\,T[p])\bigr).
\end{gather*}
La respuesta es $f(n,C)$, y si vale $+\infty$ no hay solución. Orden
$\Oh(n\,C\,\max_pQ[p])$.
\end{problema}

Con dos recursos (por ejemplo dinero y peso) la tabla tiene una dimensión
por cada uno. En los exámenes viejos la letra solía sugerir cómo definir los
casos base, distinguiendo si el requerimiento mínimo ya se alcanzó o no.\visto{Examen jul. 2009 Ej. 2: menús con dos restricciones}

\begin{problema}{Dos restricciones (menús 2009; ringtones 2010; contenedor 2014)}
Si hay dos recursos (precio $\le M$ \emph{y} proteínas $\ge W$), cada uno es
un índice. Para minimizar calorías con menús $1..k$, presupuesto $m$ y
requerimiento proteico $w$:
\begin{gather*}
f(0,m,w)=\begin{cases}0 & w\le 0\\ +\infty & w>0\end{cases}\\
f(k,m,w)=\min\bigl\{f(k-1,m,w),\ c_k+f(k-1,m-p_k,\max\{0,w-q_k\})\bigr\},
\end{gather*}
donde la segunda opción solo existe si $p_k\le m$. Truncar el requerimiento
en 0 mantiene la tabla finita: $\Oh(nMW)$. El caso base distingue si el
requerimiento mínimo ya se cumplió, que es lo que pedía la sugerencia.
\end{problema}

Hay otras variantes del programa viejo que se resuelven igual. Invertir
$d$ pesos en $n$ empresas con interés $h(x,i)$ es una mochila en la que se
elige \emph{cuánto} asignar a la última empresa:
$f(i,x)=\max_{0\le y\le x}\bigl(h(y,i)+f(i-1,x-y)\bigr)$, en
$\Oh(nd^2)$.\visto{Examen feb. 2009 Ej. 3: invertir en $n$ empresas}

Cuando el objetivo tiene dos niveles (``gastar la mayor cantidad de dinero
y, entre esas, la menor cantidad de golosinas''), el valor de cada celda es
un par que se compara de forma lexicográfica. Cuando además del máximo se
exige un mínimo (programadores con experiencia al menos $E$), basta
calcular el óptimo y verificar al final que cumple la cota.\visto{Examen feb. 2016 Ej. 4: programadores con presupuesto}

\begin{trampa}[mochila]
\begin{itemize}[leftmargin=1.2em]
  \item En la versión ilimitada, usar $\OPT(i-1,R-w_i)$: eso es 0/1.
  \item Olvidar el caso $w<w_i$, que deja índices negativos.
  \item Decir que $\Oh(nW)$ es polinomial.
  \item Proponer greedy por cociente para la versión 0/1.
  \item Omitir cómo se reconstruye el conjunto, o reconstruirlo comparando
  $M[i,w]$ con $M[i,w-1]$ en lugar de con $M[i-1,w]$.
\end{itemize}
\end{trampa}

% =====================================================================
\section{Dos secuencias: alineación, edición y subsecuencias comunes}
% =====================================================================

\subsection{Alineación de secuencias}

Dadas $X=x_1..x_m$ e $Y=y_1..y_n$, una \emph{alineación} es un conjunto $M$
de pares $(i,j)$ que forma un emparejamiento sin cruces: no hay
$(i,j),(i',j')\in M$ con $i<i'$ y $j>j'$. Cada posición sin pareja (un hueco)
cuesta $\delta$ y cada par $(i,j)\in M$ cuesta $\alpha_{x_iy_j}$ (cero si las
letras son iguales). El subproblema es sobre los dos prefijos: $\OPT(i,j)$ es
el costo mínimo de alinear $x_1..x_i$ con $y_1..y_j$.

\begin{resultado}[(6.14)]
Sea $M$ una alineación de $X$ e $Y$. Si $(m,n)\notin M$, entonces la posición
$m$ de $X$ o la posición $n$ de $Y$ no está emparejada en $M$.
\end{resultado}

\begin{resultado}[(6.15) y (6.16)]
En una alineación óptima se cumple al menos una de: $(m,n)\in M$; la
posición $m$ de $X$ no está emparejada; la posición $n$ de $Y$ no está
emparejada. Por lo tanto, para $i,j\ge1$,
\[
\OPT(i,j)=\min\bigl\{\alpha_{x_iy_j}+\OPT(i-1,j-1),\ \delta+\OPT(i-1,j),\ \delta+\OPT(i,j-1)\bigr\},
\]
con $\OPT(i,0)=i\delta$ y $\OPT(0,j)=j\delta$.
\end{resultado}

La tabla se llena por filas (o por columnas) en $\Oh(mn)$ tiempo y espacio.
El valor también es el largo del camino mínimo de $(0,0)$ a $(m,n)$ en la
grilla $G_{XY}$, con aristas horizontales y verticales de costo $\delta$ y
diagonales de costo $\alpha$ (6.17). Para el \emph{valor} alcanza con
guardar dos columnas (espacio $\Oh(m)$); para recuperar la alineación en
espacio lineal está el algoritmo de Hirschberg, que combina esa idea con
dividir y conquistar en tiempo $\Oh(mn)$. Saber ubicar el problema (qué
resuelve y con qué técnica) se preguntó como teoría.\visto{Examen feb. 2018 Preg. 4: qué técnica usa alineamiento}

\subsection{Distancia de edición}

Es la alineación con $\delta=1$ y $\alpha_{pq}=1$ si $p\ne q$ (0 si son
iguales): un par desigual es una sustitución, una posición de $x$ sin pareja
es un borrado y una de $y$ sin pareja es una inserción.\visto{Examen feb. 2019 Ej. 2: distancia de edición}

\begin{problema}{Distancia de edición (examen feb. 2019, Ej. 2)}
\textbf{Base.} $\OPT(0,j)=j$ (hay que insertar $j$ caracteres, y no se puede
con menos) y $\OPT(i,0)=i$ (borrar $i$).

\textbf{Lema que se pedía.} En un emparejamiento óptimo $S$, la posición $n$
de $x$ o la $m$ de $y$ está emparejada: si ninguna lo estuviera, agregar
$(n,m)$ cambia un borrado y una inserción por una sustitución y baja el
costo. Si ambas lo están, es entre sí, porque si no habría un cruce.

\textbf{Recurrencia.} Con $\delta_{i,j}=[x_i\ne y_j]$,
\[
\OPT(i,j)=\min\bigl\{\OPT(i-1,j-1)+\delta_{i,j},\ \OPT(i-1,j)+1,\ \OPT(i,j-1)+1\bigr\},\quad i,j>0.
\]
Cada término corresponde a: $(i,j)\in S$; $x_i$ se borra; $y_j$ se inserta.
En los tres casos lo que queda de $S$ debe ser óptimo para los prefijos
correspondientes (corte y pegado). Orden $\Oh(nm)$.
\end{problema}

En diciembre de 2018 apareció el mismo molde con otra historia: dos programas
que se ejecutan a la vez y comparten un ciclo cuando la próxima instrucción
coincide. La letra preguntaba primero por los casos base, que es la forma
habitual de arrancar estos ejercicios.\visto{Examen dic. 2018 Ej. 3: ejecución paralela de dos programas}

\begin{problema}{Ejecución en paralelo de dos programas (examen dic. 2018, Ej. 3)}
En cada ciclo, si las próximas instrucciones de $X$ e $Y$ son iguales se
ejecutan juntas; si no, se ejecuta una sola. Con $\OPT(i,j)$ los ciclos
mínimos para las primeras $i$ de $X$ y $j$ de $Y$:
\begin{gather*}
\OPT(0,j)=j,\qquad \OPT(i,0)=i,\\
\OPT(i,j)=\begin{cases}1+\OPT(i-1,j-1) & x_i=y_j\\ 1+\min\{\OPT(i-1,j),\OPT(i,j-1)\} & x_i\ne y_j.\end{cases}
\end{gather*}
Cuando $x_i=y_j$ conviene ejecutarlas juntas: cualquier ejecución óptima en
la que no lo están se transforma en otra igual de buena postergando la
primera al último ciclo. Observación: el resultado es $n+m-\mathrm{LCS}(X,Y)$.
\end{problema}

\subsection{Subsecuencia y substring común más largos}

La subsecuencia común más larga (LCS) es el mismo molde con máximo:
$\OPT(i,j)=1+\OPT(i-1,j-1)$ si $x_i=y_j$, y
$\max\{\OPT(i-1,j),\OPT(i,j-1)\}$ si no, con bordes en 0. El parcial de 2017
daba el código de la variante para \emph{substrings} (caracteres
consecutivos) y pedía descifrarlo.\visto{2P 2017 Ej. 1: substring común más largo}

\begin{problema}{Substring común más largo (2P 2017, Ej. 1)}
El código usa dos tablas. $\mathit{Mp}[i,j]$ es el largo del sufijo común más largo
de $A_1..A_i$ y $B_1..B_j$ (el substring común que \emph{termina} en $i$ y
en $j$); $\mathit{Mt}[i,j]$ es el largo del substring común más largo de esos
prefijos:
\begin{gather*}
\mathit{Mp}[i,j]=\begin{cases}1+\mathit{Mp}[i-1,j-1] & A_i=B_j\\ 0 & A_i\ne B_j\end{cases}\\
\mathit{Mt}[i,j]=\max\{\mathit{Mp}[i,j],\mathit{Mt}[i-1,j],\mathit{Mt}[i,j-1]\},
\end{gather*}
con fila y columna 0 en cero. Para $A=\texttt{la}$, $B=\texttt{ala}$:

\begin{center}
\begin{tabular}{c|ccc}
$\mathit{Mp}$ & a & l & a\\\hline
l & 0 & 1 & 0\\
a & 1 & 0 & 2
\end{tabular}\qquad
\begin{tabular}{c|ccc}
$\mathit{Mt}$ & a & l & a\\\hline
l & 0 & 1 & 1\\
a & 1 & 1 & 2
\end{tabular}
\end{center}

El resultado es 2 (\texttt{la}). La versión recursiva sin tabla recalcula
los mismos $\mathit{Mp}(i,j)$ muchas veces. El espacio $\Oh(nm)$ baja a $\Oh(m)$
guardando solo la fila anterior.
\end{problema}

En 2013 se pidió la alineación de dos secuencias de cubos de colores con un
cubo ``comodín'' para los huecos. Es exactamente la sección 6.6 con la
función de costos de la letra.\visto{Examen jul. 2013 Prob. 1: alinear cubos de colores}

\begin{receta}[dos secuencias]
\begin{itemize}[leftmargin=1.2em]
  \item Definí $\OPT(i,j)$ sobre \emph{prefijos} $x_1..x_i$, $y_1..y_j$ con
  $0\le i\le n$, $0\le j\le m$, y escribí la fila y la columna 0.
  \item Justificá con ``¿qué pasa con las últimas posiciones $i$ y $j$?'': se
  emparejan entre sí, $x_i$ queda solo o $y_j$ queda solo.
  \item Llená con dos \textbf{para} anidados crecientes y devolvé $\OPT(n,m)$.
  \item Para reconstruir, partí de $(n,m)$ y en cada celda andá a la vecina
  (diagonal, arriba o izquierda) cuyo término dio el mínimo.
\end{itemize}
\end{receta}

% =====================================================================
\section{\texorpdfstring{Intervalos $i..j$}{Intervalos i..j}}\label{pd:sec:intervalos}
% =====================================================================

Cuando una decisión parte el problema en un pedazo izquierdo y uno derecho
que \emph{no} son prefijos, los subproblemas tienen que ser intervalos
arbitrarios $x_i..x_j$. Hay $\Theta(n^2)$ de ellos, y la tabla se llena por
largo creciente $\ell=j-i$ (o con $i$ decreciente y $j$ creciente), para que
los intervalos más cortos estén calculados.

\subsection{Estructura secundaria del ARN}

Una secuencia $B=b_1..b_n$ sobre $\{A,C,G,U\}$; una estructura secundaria es
un conjunto $S$ de pares $(i,j)$ tal que: no hay giros cerrados ($i<j-4$);
cada par es $\{A,U\}$ o $\{C,G\}$; ninguna base está en dos pares; y no hay
cruces (no existen $(i,j),(k,l)\in S$ con $i<k<j<l$). Se busca $|S|$
máximo.

Con $\OPT(j)$ sobre prefijos no alcanza. Si $b_j$ se empareja con $b_t$, por
la condición de no cruce el problema se parte en $b_1..b_{t-1}$ y
$b_{t+1}..b_{j-1}$, y el segundo no es un prefijo. Por eso se usa
$\OPT(i,j)$ sobre $b_i..b_j$.

\begin{resultado}[Recurrencia del ARN (sección 6.5)]
$\OPT(i,j)=0$ si $i\ge j-4$. Si no,
\[
\OPT(i,j)=\max\Bigl(\OPT(i,j-1),\ \max_{t}\bigl(1+\OPT(i,t-1)+\OPT(t+1,j-1)\bigr)\Bigr),
\]
donde $t$ recorre los $i\le t<j-4$ tales que $b_t$ y $b_j$ forman un par
permitido. Hay $\Oh(n^2)$ subproblemas de costo $\Oh(n)$: $\Oh(n^3)$.
\end{resultado}

\pdalgo{Llenado por largo de intervalo (ARN)}
\begin{algorithmic}[1]
\State $M[i,j]\gets0$ para todo $i\ge j-4$
\For{$k=5$ \textbf{hasta} $n-1$} \Comment{$k=j-i$}
  \For{$i=1$ \textbf{hasta} $n-k$}
    \State $j\gets i+k$
    \State $M[i,j]\gets\max\bigl(M[i,j-1],\ \max_t(1+M[i,t-1]+M[t+1,j-1])\bigr)$
  \EndFor
\EndFor
\State \Return $M[1,n]$
\end{algorithmic}

\subsection{Multiplicación de una cadena de matrices}

No está en el capítulo 6 del libro, pero es el ejemplo más directo del
molde: la última multiplicación parte la cadena en dos
subcadenas, y cada una se resuelve de forma óptima.\visto{Examen jul. 2021 Ej. 2: producto de cadena de matrices}

\begin{problema}{Producto $A_1\times\dots\times A_n$ (examen jul. 2021, Ej. 2)}
$A_i$ es de $d_{i-1}\times d_i$. Con $\OPT(i,j)$ el mínimo de
multiplicaciones escalares para $A_i\times\dots\times A_j$:
\[
\OPT(i,i)=0,\qquad \OPT(i,j)=\min_{i\le k<j}\bigl(\OPT(i,k)+\OPT(k+1,j)+d_{i-1}d_kd_j\bigr),\quad i<j.
\]
Calcular $A_i..A_j$ es calcular $U=A_i..A_k$, $V=A_{k+1}..A_j$ (de forma
óptima, por definición de $\OPT$) y $U\times V$, que cuesta
$d_{i-1}d_kd_j$. Se llena con $\ell=1..n-1$ e $i=1..n-\ell$, $j=i+\ell$;
respuesta $\OPT(1,n)$, orden $\Oh(n^3)$.
\end{problema}

Con la misma estructura se tomó en 2011 ``intercalar paréntesis en una suma
para minimizar la suma de las sumas intermedias'': el término de combinación
es $\sum_{t=i}^{j}x_t$ (que se obtiene en $\Oh(1)$ con sumas
acumuladas) en lugar de $d_{i-1}d_kd_j$.\visto{Examen dic. 2011 Prob. 2: paréntesis con sumas intermedias}

\subsection{Palíndromos}

\begin{problema}{Subsecuencia palíndroma más larga (examen jul. 2022, Ej. 2)}
Con $\OPT(i,j)$ el largo máximo de una subsecuencia de $x_i..x_j$ que es
palíndromo:
\begin{gather*}
\OPT(i,i)=1,\qquad \OPT(i,j)=0\quad (j<i),\\
\OPT(i,j)=\begin{cases}2+\OPT(i+1,j-1) & x_i=x_j\\ \max\{\OPT(i+1,j),\OPT(i,j-1)\} & x_i\ne x_j\end{cases}\quad (i<j).
\end{gather*}
Si $x_i=x_j$, hay un palíndromo óptimo que usa ambos extremos (si usa solo
uno, se lo cambia por el otro; si no usa ninguno, se lo puede alargar). Si
$x_i\ne x_j$, no puede usar los dos. Orden de llenado: $i$ desde $n-1$
bajando hasta 1 y $j$ desde $i+1$ subiendo hasta $n$. $\Oh(n^2)$.
\end{problema}

Para \emph{substrings} palíndromos (lo que necesita la segmentación de 2026)
la tabla booleana es $\mathit{esP}(i,i)=\texttt{true}$,
$\mathit{esP}(i,i+1)=(a_i=a_{i+1})$ y $\mathit{esP}(i,j)=(a_i=a_j)\wedge \mathit{esP}(i+1,j-1)$, con el
mismo orden de llenado; se calcula en $\Oh(n^2)$.\visto{Examen jul. 2022 Ej. 2: subsecuencia palíndroma más larga}

\begin{trampa}[intervalos]
\begin{itemize}[leftmargin=1.2em]
  \item Llenar con $i$ creciente: $\OPT(i,j)$ necesita $\OPT(i+1,\cdot)$, que
  todavía no está.
  \item Olvidar los intervalos vacíos ($j<i$) cuando la recurrencia los usa.
  \item Plantear $\OPT(j)$ sobre prefijos cuando la decisión deja un
  pedazo ``del medio''.
\end{itemize}
\end{trampa}

% =====================================================================
\section{Grillas y tableros}
% =====================================================================

Una grilla en la que solo se avanza en ciertas direcciones (hacia arriba,
hacia abajo, a la derecha) es un DAG implícito: cada casilla es un
subproblema y depende de sus dos o tres vecinas ``anteriores''. La tabla es
la grilla misma y se llena fila por fila en el sentido del movimiento. Los
bordes se manejan con centinelas $\pm\infty$ en columnas ficticias, que es lo
que sugieren las letras.\visto{2P 2023 Ej. 1: sargento en un tablero}

\begin{problema}{El sargento (2P 2023, Ej. 1)}
Desde $(i,j)$ el sargento va a $(i-1,j-1)$, $(i-1,j)$ o $(i-1,j+1)$; salir
de $(i,j)$ cuesta $C(i,j)$. Se va de $(n,1)$ a alguna casilla de la fila 1.
Con $\OPT(i,j)$ el costo mínimo de llegar a $(i,j)$:
\[
\OPT(n,1)=0,\quad \OPT(n,j)=\infty\ (j>1),\quad \OPT(i,0)=\OPT(i,n+1)=\infty,
\]
\[
\OPT(i,j)=\min_{\ell\in\{j-1,j,j+1\}}\bigl(\OPT(i+1,\ell)+C(i+1,\ell)\bigr),\quad 1\le i<n,\ 1\le j\le n.
\]
Se llena con $i$ desde $n-1$ bajando hasta 1; la respuesta es
$\min_k\OPT(1,k)$. Para reconstruir, se parte de la columna que da el mínimo
en la fila 1 y en cada fila se baja a la vecina cuyo término coincide con el
valor actual. $\Oh(n^2)$.
\end{problema}

El minero de 2022 es el espejo: la máquina perfora hacia abajo desde una
columna inicial a elección, maximizando el valor extraído, con
$v_{i,0}=v_{i,n+1}=-\infty$ como centinelas.\visto{Examen feb. 2022 Ej. 3: perforación de una veta}

\begin{problema}{Perforación minera (examen feb. 2022, Ej. 3)}
Con $\OPT(i,j)$ la ganancia máxima perforando desde $(i,j)$ (incluida) hasta
la profundidad $k$:
\begin{gather*}
\OPT(k,j)=v_{k,j},\\
\OPT(i,j)=v_{i,j}+\max\{\OPT(i+1,j-1),\OPT(i+1,j),\OPT(i+1,j+1)\},\quad 1\le i<k.
\end{gather*}
Respuesta $\max_j\OPT(1,j)$; el plan es la columna del $\pdargmax$ y, en cada
fila, la dirección $I$, $V$ o $D$ cuyo término da el valor de la celda.
$\Oh(kn)$.
\end{problema}

\begin{problema}{Contar caminos y recolectar en una grilla (jul. 2017; jul. 2018)}
\textbf{Contar.} Caminos de $(i,j)$ a $(m,n)$ moviéndose a la derecha o
hacia arriba: $\mathit{NC}(i,j)=1$ si $i=m$ o $j=n$, y
$\mathit{NC}(i,j)=\mathit{NC}(i+1,j)+\mathit{NC}(i,j+1)$ en otro caso. Los caminos se parten en
disjuntos según el primer paso, por eso se \emph{suma}. Para $m=6$, $n=3$
da $\binom{9}{3}=84$.

\textbf{Recolectar.} Con $E(i,j)$ objetos en la esquina $(i,j)$ y calles
flechadas hacia índices crecientes: $C(i,j)=0$ si $i=0$ o $j=0$, y
$C(i,j)=\max\{C(i-1,j),C(i,j-1)\}+E(i,j)$. Si hay esquinas intransitables
($E=-1$), se define la accesibilidad $A(i,j)$ (la esquina es transitable y
alguna de las dos anteriores es accesible) y se toma $C(i,j)=-\infty$ en las
no accesibles.
\end{problema}

Estos dos son del programa viejo, pero el de contar caminos marca un
patrón que volvió en 2019: cuando la pregunta es ``¿cuántos?'', los casos
se \emph{suman} en lugar de compararse, porque particionan el conjunto
de soluciones.\visto{Examen jul. 2017 Prob. A: caminos en una grilla}

El del recolector agrega una idea reutilizable: cuando algunas celdas están
prohibidas, conviene una recurrencia auxiliar booleana de alcanzabilidad (o
el valor $-\infty$) para no confundir ``no se puede llegar'' con ``se llega
sin recolectar nada''.\visto{Examen jul. 2018 Prob. A: recolector en calles flechadas}

\begin{trampa}[grillas]
Olvidar los bordes (consultar la columna 0 o $n+1$ sin haberlas definido);
devolver $\OPT(1,1)$ cuando el destino es ``cualquier casilla de la fila
superior''; recorrer las filas en el sentido opuesto al movimiento.
\end{trampa}

% =====================================================================
\section{DAGs, conteo y árboles}
% =====================================================================

\subsection{PD sobre un DAG}

En un grafo dirigido acíclico, cualquier recurrencia en la que $f(v)$ se
obtiene combinando los valores $f(w)$ de los sucesores $w$ de $v$ (con un
máximo, un mínimo o una suma) se calcula procesando
los vértices en orden topológico inverso: cuando se llega a $v$, todos sus
sucesores ya están resueltos. Cada arista se mira una vez, así que el tiempo
es $\Oh(n+m)$ con listas de adyacencia. El orden topológico se obtiene con
el algoritmo de los grados de entrada del capítulo 3, que se pidió
reescribir en diciembre de 2018 y en julio de 2026.\visto{Examen feb. 2023 Ej. 3: camino más largo en un DAG}

\begin{problema}{Camino más largo en un DAG (examen feb. 2023, Ej. 3)}
Los vértices $1..n$ ya están en orden topológico. Con $L(v)$ la longitud
(en aristas) del camino más largo que empieza en $v$:
\[
L(v)=\begin{cases}0 & v\text{ no tiene aristas salientes}\\ 1+\max_{(v,w)\in E}L(w) & \text{si no.}\end{cases}
\]
Se calcula con $v$ desde $n$ bajando hasta 1 y la respuesta es $\max_vL(v)$.
Para devolver el camino, guardar $\mathit{sig}[v]$ (el $w$ que dio el máximo) y el
vértice de inicio del mejor, y seguir los punteros. $\Oh(|V|+|E|)$.
\end{problema}

Cuando el orden topológico no viene dado, se calcula dentro del mismo
algoritmo. En julio de 2026 los valores se propagaban hacia adelante mientras
se sacaban vértices de la cola, y había que detectar además si el grafo tenía
ciclos.\visto{Examen jul. 2026 Ej. 2: tiempos de inicio con dependencias}

\begin{problema}{Tiempos de inicio con dependencias (examen jul. 2026, Ej. 2)}
La subtarea $i$ dura $d_i$ y solo puede empezar cuando terminaron todas las
que la preceden. El inicio mínimo es $s_i=0$ si no tiene previas y
$s_i=\max\{s_j+d_j:\ i\in p_j\}$ si no (camino crítico). Se calcula en el
orden en que el algoritmo de grados de entrada saca los vértices de la cola:
al sacar $h$, para cada $i\in p_h$ se hace $s_i\gets\max\{s_i,s_h+d_h\}$, se
decrementa $a_i$ y se encola $i$ si llegó a 0. Si al final queda algún
$a_i>0$ hay un ciclo y el proyecto no es factible. $\Oh(n+m)$.

\textbf{Error típico.} Recorrer $i=1..n$ en orden de índice y propagar
$s_j\gets\max\{s_j,s_i+d_i\}$: puede usar un $s_i$ todavía no definitivo.
Con $p_1=\{\}$, $p_2=\{1\}$, $p_3=\{2\}$ y $d=(5,10,15)$ da $s_1=10$ en vez
de 25.
\end{problema}

\subsection{Contar caminos}

Para contar caminos de $v$ a $t$ en un DAG: $N(t)=1$ y
$N(v)=\sum_{(v,w)\in E}N(w)$. Para contar caminos \emph{de largo mínimo}
desde un origen se restringe la suma a las aristas que respetan las
distancias: en un grafo sin pesos, las que van de un nivel BFS al
siguiente.\visto{Examen jul. 2019 Ej. 2: contar historias más cortas}

\begin{problema}{Historias más cortas de un libro (examen jul. 2019, Ej. 2)}
Los capítulos forman un DAG con a lo sumo 3 aristas por vértice; se pide el
largo mínimo $L$ de una historia (de 1 a un capítulo final) y la cantidad
$N$ de historias de ese largo, en $\Oh(n)$. Con $L(w)$ el nivel BFS de $w$
desde 1, los caminos mínimos a $w$ se parten según el último vértice antes
de $w$, que está en el nivel anterior:
\[
N(1)=1,\qquad N(w)=\sum_{t:\ (t,w)\in E,\ L(w)=L(t)+1}N(t).
\]
Se hace BFS por niveles acumulando $N$ y se corta al completar el primer
nivel con capítulos finales: $L$ es ese nivel y $N$ es la suma de $N(u)$
sobre los finales de ese nivel. El análisis es el de BFS con $|E|=\Oh(n)$.
\end{problema}

\subsection{Árboles}

En un árbol con raíz, los subproblemas son los subárboles y se resuelven en
postorden (hijos antes que el padre). El caso clásico es el conjunto
independiente de peso máximo (capítulo 10 del libro): para cada vértice $u$
se calculan dos valores, el óptimo del subárbol de $u$ que incluye a $u$ y
el que no lo incluye.\visto{Examen feb. 2022 Ej. 1: independiente de peso máximo en árbol}

\begin{problema}{Conjunto independiente de peso máximo en árbol binario (examen feb. 2022, Ej. 1)}
$q(T)$ es el peso máximo de un independiente de $T$ que no incluye la raíz
$v$; $p(T)$ el máximo sin restricción. Si $T$ es una hoja, $p(T)=w(v)$ y
$q(T)=0$. Si no, con subárboles $T_I$ y $T_D$:
\[
q(T)=p(T_I)+p(T_D),\qquad p(T)=\max\bigl\{w(v)+q(T_I)+q(T_D),\ q(T)\bigr\}.
\]
La unión de independientes de $T_I$ y $T_D$ es independiente; si se incluye
$v$ no se pueden incluir las raíces de los subárboles. La función recursiva
devuelve el par $(p,q)$; para un árbol perfecto,
$f(n)\le2f\bigl(\tfrac{n-1}{2}\bigr)+c$, que da $f(n)\le cn$ por inducción.
En general es $\Oh(n)$ porque cada nodo se procesa una vez.
\end{problema}

Sin pesos, en un bosque alcanza con el greedy que toma hojas (diciembre
de 2019, demostrado por intercambio). Con pesos ese greedy falla: en el
camino $a-b-c$ con pesos $1,10,1$ toma $\{a,c\}$ (peso 2) y el óptimo es
$\{b\}$ (peso 10). Por eso con pesos hace falta la PD.

\begin{receta}[DAGs y árboles]
\begin{itemize}[leftmargin=1.2em]
  \item Decí explícitamente en qué orden se procesan los vértices (orden
  topológico inverso, postorden) y por qué así los valores necesarios ya
  están calculados.
  \item Para $\Oh(n+m)$: cada vértice se procesa una vez y cada arista se mira
  una vez desde su origen, con listas de adyacencia.
  \item Si se cuenta, sumá sobre una partición de las soluciones y decí cuál
  es la partición.
\end{itemize}
\end{receta}

% =====================================================================
\section{Caminos mínimos con costos negativos: Bellman--Ford}
% =====================================================================

\subsection{Por qué no alcanza Dijkstra}

Con costos negativos, Dijkstra puede fijar mal una distancia: con aristas
$s\to a$ de costo 2, $s\to b$ de costo 3 y $b\to a$ de costo $-2$, fija
$d(a)=2$ cuando el camino $s,b,a$ cuesta 1. Sumarle una constante a todas las
aristas para hacerlas positivas tampoco sirve, porque penaliza más a los
caminos con más aristas. Y si hay un ciclo de costo negativo desde el que se
alcanza $t$, el costo mínimo no está definido (se puede dar vueltas
indefinidamente).

\begin{resultado}[(6.22)]
Si $G$ no tiene ciclos negativos, existe un camino mínimo de $s$ a $t$ que es
simple y, por lo tanto, tiene a lo sumo $n-1$ aristas.
\end{resultado}

Esto sugiere el parámetro: la \emph{cantidad de aristas permitidas}. En el
libro los caminos van \emph{hacia} $t$: $\OPT(i,v)$ es el costo mínimo de un
camino de $v$ a $t$ con a lo sumo $i$ aristas, con $\OPT(0,t)=0$ y
$\OPT(0,v)=\infty$ para $v\ne t$. La respuesta es $\OPT(n-1,s)$.

\begin{resultado}[(6.23)]
Si $i>0$, entonces
$\OPT(i,v)=\min\Bigl(\OPT(i-1,v),\ \min_{w\in V}\bigl(\OPT(i-1,w)+c_{vw}\bigr)\Bigr)$.
\end{resultado}

O el camino óptimo usa a lo sumo $i-1$ aristas (primer término), o usa
exactamente $i$ y la primera es alguna $(v,w)$; lo que sigue es un camino de
$w$ a $t$ con a lo sumo $i-1$ aristas, que tiene que ser óptimo.

\pdalgo{Bellman--Ford (Shortest-Path$(G,s,t)$)}
\begin{algorithmic}[1]
\State $M[0,t]\gets0$;\ $M[0,v]\gets\infty$ para todo $v\ne t$
\For{$i=1$ \textbf{hasta} $n-1$}
  \ForAll{$v\in V$}
    \State $M[i,v]\gets\min\bigl(M[i-1,v],\ \min_{(v,w)\in E}(M[i-1,w]+c_{vw})\bigr)$
  \EndFor
\EndFor
\State \Return $M[n-1,s]$
\end{algorithmic}

\begin{resultado}[(6.24) y (6.25)]
Si $G$ no tiene ciclos negativos, el algoritmo calcula correctamente el
costo mínimo de un camino $s$-$t$. La tabla tiene $n^2$ celdas y cada una
cuesta $\Oh(n)$, lo que da $\Oh(n^3)$; mirando solo las aristas que salen de
$v$, cada iteración cuesta $\Oh(\sum_vn_v)=\Oh(m)$ y el total es $\Oh(mn)$.
\end{resultado}

En la práctica se usa un solo arreglo $M[v]$ actualizado en el lugar, más un
puntero $\mathit{first}[v]$ al primer vértice del camino actual, lo que baja el
espacio a $\Oh(n)$ y permite cortar en cuanto una iteración no cambia nada.
Si el grafo de punteros $\{(v,\mathit{first}[v])\}$ llega a tener un ciclo, ese ciclo
tiene costo negativo.

\subsection{Las versiones de examen}

En el parcial de 2018 se pidió primero la recurrencia para ``exactamente
$i$ aristas'' y después obtener el óptimo con a lo sumo $K$ aristas como un
mínimo; la letra fijaba los nombres y había que respetarlos.\visto{2P 2018 Ej. 2: caminos con a lo sumo $K$ aristas}

\begin{problema}{Caminos mínimos con a lo sumo $K$ aristas (2P 2018, Ej. 2)}
$\mathit{CMC}(i,v)$ es el costo mínimo de un camino de $v$ a $t$ con
\emph{exactamente} $i$ aristas ($\infty$ si no hay):
\[
\mathit{CMC}(0,v)=\begin{cases}0&v=t\\ \infty&v\ne t\end{cases}\qquad
\mathit{CMC}(i,v)=\min_{(v,w)\in E}\bigl(\mathit{CMC}(i-1,w)+c_{vw}\bigr),\quad 1\le i\le K,
\]
con $\min\emptyset=\infty$, y
$\OPT(K,v)=\min_{0\le i\le K}\mathit{CMC}(i,v)$. La alternativa directa es
$\OPT(k,t)=0$ y $\OPT(k,v)=\min_{(v,w)}(\OPT(k-1,w)+c_{vw})$ para $v\ne t$.
Con listas de adyacencia el ciclo principal es $\Oh(K(n+m))$ (con matriz,
$\Oh(Kn^2)$), y el mínimo final $\Oh(nK)$; como $K<n$, es polinomial.
\end{problema}

En julio de 2023 se pidió lo mismo pero con caminos que salen del origen, y
con el costo de cada arista armado a partir de dos datos de la letra. Los
costos son positivos, así que Dijkstra también serviría; lo que se evaluaba
era la recurrencia en $i$.\visto{Examen jul. 2023 Ej. 2: routers con a lo sumo $i$ saltos}

\begin{problema}{Routers con a lo sumo $i$ saltos (examen jul. 2023, Ej. 2)}
Cada enlace $\{u,v\}$ da dos aristas $(u,v)$ y $(v,u)$; la arista $(u,v)$
cuesta $\mathit{TP}(u)+\mathit{TT}(u,v)$ (procesar en $u$ y transmitir). Aquí los caminos
salen \emph{desde} $s$: $\OPT(0,s)=0$, $\OPT(0,v)=\infty$ para $v\ne s$ y
\begin{multline*}
\OPT(i,v)=\min\Bigl\{\OPT(i-1,v),\\ \min_{(u,v)\in E}\bigl(\OPT(i-1,u)+\mathit{TP}(u)+\mathit{TT}(u,v)\bigr)\Bigr\},\quad 1\le i\le n-1.
\end{multline*}
Respuesta $\OPT(n-1,t)$, en $\Oh(nm)$. Para el camino: desde $v=t$ e $i=n-1$
bajando, si $\OPT(i,v)\ne\OPT(i-1,v)$ se busca el $u$ que dio el mínimo, se
lo agrega al principio y $v\gets u$.
\end{problema}

\subsection{Detección de ciclos negativos}

Para saber si $G$ tiene \emph{algún} ciclo negativo (no solo uno que llegue a
$t$) se agrega un vértice nuevo $t$ y una arista $(v,t)$ de costo 0 desde
cada $v$: todo ciclo de $G$ queda conectado a $t$, y $G$ tiene un ciclo
negativo si y solo si el grafo aumentado tiene uno desde el que se llega a
$t$. Después se corre Bellman--Ford una iteración más de lo necesario.

\begin{resultado}[(6.30) y (6.31)]
Si $\OPT(n,v)=\OPT(n-1,v)$ para todo $v$, no hay ciclos negativos que
lleguen a $t$. Si $\OPT(n,v)<\OPT(n-1,v)$ para algún $v$, un camino de $v$ a
$t$ de costo $\OPT(n,v)$ tiene un ciclo, y ese ciclo tiene costo negativo.
El ciclo se encuentra siguiendo ese camino, y todo se hace en $\Oh(mn)$.
\end{resultado}

La aplicación típica es el arbitraje: con tasas de cambio $r_{ij}$ entre
monedas, hay un ciclo de cambios con ganancia ($\prod r_{ij}>1$) si y solo si
hay un ciclo negativo con costos $c_{ij}=-\log r_{ij}$. En los exámenes
recientes no se tomó la detección de ciclos, pero es teoría del capítulo
que se puede preguntar como verdadero o falso.

\subsection{Floyd--Warshall (programa anterior)}

Hasta 2017 el curso daba Floyd para caminos mínimos entre todo par. No está
en el capítulo 6 del libro, pero los exámenes viejos lo preguntan seguido y
la idea es la misma: el parámetro extra limita los vértices intermedios
permitidos.\visto{Examen jul. 2015 Ej. 4: optimalidad y recurrencia de Floyd}

\begin{resultado}[Floyd]
Con $f_k(i,j)$ el costo mínimo de un camino de $i$ a $j$ cuyos vértices
intermedios están en $\{1..k\}$: $f_0(i,j)=W(i,j)$ y
$f_k(i,j)=\min\{f_{k-1}(i,j),\ f_{k-1}(i,k)+f_{k-1}(k,j)\}$. Se puede usar
una sola matriz; el tiempo es $\Theta(n^3)$. Funciona con costos negativos si
no hay ciclos negativos, y un valor negativo en la diagonal indica uno.
\end{resultado}

El principio de optimalidad que se pedía enunciar es que todo subcamino de
un camino mínimo es mínimo (si no, se lo reemplaza). Por eso el camino
mínimo de $v_1$ a $v_k$ que visita $v_2..v_{k-1}$ en orden es la
concatenación de los caminos mínimos entre consecutivos.\visto{2P 2009 Ej. 3: camino que visita una secuencia}

La variante de carga máxima de 2011 (el ``ancho'' de un camino es su arista
de menor capacidad y se quiere el camino más ancho) cambia las operaciones:
\[
f_k(i,j)=\max\bigl\{f_{k-1}(i,j),\ \min\{f_{k-1}(i,k),f_{k-1}(k,j)\}\bigr\}.
\]
Para reconstruir se guarda en $\mathit{Camino}[i][j]$ el último $k$ que mejoró la
celda y se arma recursivamente $i\leadsto k\leadsto j$.\visto{2P 2011 Ej. 2: carga máxima entre ciudades}

\begin{trampa}[caminos mínimos]
\begin{itemize}[leftmargin=1.2em]
  \item Usar Dijkstra (o sumar una constante) con costos negativos.
  \item Olvidar el término $\OPT(i-1,v)$ en la versión ``a lo sumo $i$
  aristas'' (sin él, es ``exactamente $i$'').
  \item Confundir el sentido: en el libro $\OPT(i,v)$ es de $v$ \emph{hacia}
  $t$ y se minimiza sobre aristas que \emph{salen} de $v$; si se define desde
  $s$, se minimiza sobre aristas que \emph{entran} a $v$.
  \item Analizar como $\Oh(n^3)$ cuando se pide $\Oh(mn)$: hay que decir que
  cada arista se mira una vez por iteración.
  \item Afirmar que Bellman--Ford ``funciona con ciclos negativos'': los
  detecta, pero con ellos no hay camino mínimo.
\end{itemize}
\end{trampa}

% =====================================================================
\section{Mapa de lo que se preguntó}
% =====================================================================

Todos los ejercicios de programación dinámica de los parciales y exámenes
disponibles, por año. Hasta 2017 el curso seguía otro programa (con
backtracking y Floyd, y el pedido típico era ``dé una fórmula recursiva
$f$''); desde 2018 sigue el libro de Kleinberg y Tardos y el pedido es
recurrencia, algoritmo iterativo, reconstrucción y orden.

\newcommand{\pdmapa}[1]{%
\par\smallskip{\small
\begin{tabularx}{\textwidth}{@{}>{\raggedright\arraybackslash}p{2.9cm}>{\raggedright\arraybackslash}p{2.5cm}>{\raggedright\arraybackslash}X@{}}
\toprule
\textbf{Fuente} & \textbf{Patrón} & \textbf{Qué se pedía}\\
\midrule
#1
\bottomrule
\end{tabularx}}\par\smallskip}

\pdmapa{
2P 2008 Ej.~1 & Mochila 0/1 & Fórmula recursiva para animales en un barco de capacidad $G$.\\
Ex. feb. 2009 Ej.~3 & Capacidad (asignación) & Interés óptimo invirtiendo $d$ pesos en $n$ empresas; versión iterativa y cómo reconstruir.\\
Ex. jul. 2009 Ej.~2 & Dos restricciones & Menús de mínimo valor calórico con precio $\le M$ y proteínas $\ge W$.\\
2P 2009 Ej.~3 & Caminos mínimos & Camino mínimo de $v_1$ a $v_k$ que visita la secuencia en orden.\\
Ex. dic. 2009 Prob.~1 & Capacidad con stock & Golosinas: gastar lo más posible con la menor cantidad.\\
Ex. jul. 2010 Prob.~2 & Dos restricciones & Ringtones de tamaño mínimo con presupuesto y ranking mínimo.\\
2P 2010 Ej.~2 & Prefijo con estado & Dos integrantes marcan lugares en orden; $f(l,k)$.\\
Ex. dic. 2010 Prob.~1 & Mochila 0/1 & Comerciante romano: contraejemplo greedy y recurrencia.\\
2P 2011 Ej.~2 & Floyd (máx-mín) & Carga máxima entre todo par de ciudades; implementación y camino.\\
Ex. dic. 2011 Prob.~2 & Intervalos & Paréntesis que minimizan la suma de sumas intermedias.\\
Ex. feb. 2012 Ej.~2 & Mochila 0/1 & Publicidad con presupuesto: greedy falla; recurrencia del rating.\\
Ex. jul. 2012 Ej.~2 & Floyd & Qué resuelve, qué técnica usa y aplicación paso a paso.\\
2P 2012 Ej.~1 & Capacidad exacta con stock & Carga exacta $C$ con mínimas unidades; $f(p,c)$, iterativo y orden.\\
Ex. feb. 2013 Ej.~2 & Teoría & Enunciar el principio de optimalidad y aplicarlo a caminos mínimos.\\
Ex. jul. 2013 Ej.~1 & Memoización & Verdadero o falso: Fibonacci se puede calcular en $\Oh(n)$.\\
Ex. jul. 2013 Prob.~1 & Dos secuencias & Alineación óptima de cubos de colores con comodines.\\
}

\pdmapa{
Ex. jul. 2014 Ej.~3 & Dos restricciones & Contenedor: máximo de compras con capacidad $C$ y valor $\le M$.\\
Ex. jul. 2015 Ej.~4 & Floyd & Por qué vale el principio de optimalidad; recurrencia de Floyd.\\
2P 2015 Ej.~1 & Prefijo, $p(i)=i-d$ & Edificios a distancia $\ge d$; greedy falla; recurrencia e implementación.\\
Ex. dic. 2015 Ej.~2 & Mochila 0/1 & Definir la variante, demostrar optimalidad, plantear la recurrencia.\\
Ex. feb. 2016 Ej.~4 & Dos restricciones & Programadores: experiencia máxima con presupuesto y mínimo $E$.\\
Ex. jul. 2016 Ej.~4 & Teoría & Qué resuelven Floyd, suma de subconjuntos y otros, y con qué técnica.\\
Ex. jul. 2016 Prob.~C & Mochila 0/1 & Técnica, implementación en C* y versión fraccionaria.\\
2P 2016 Ej.~2 & Último elemento & Largos $L_i$ de subsecuencias crecientes; versión recursiva y PD.\\
Ex. feb. 2017 Preg.~3 & Memoización & Fibonacci con DyC y con PD en espacio $\Oh(1)$; contar sumas.\\
Ex. jul. 2017 Prob.~A & Grilla (conteo) & Cantidad de caminos monótonos $\mathit{NC}(i,j)$; caso $m=6$, $n=3$.\\
2P 2017 Ej.~1 & Dos secuencias & Descifrar el código del substring común más largo; memoria.\\
Ex. dic. 2017 Prob.~B & Prefijo binario & Baldosas no contiguas de máximo valor; devolver las elegidas.\\
Ex. feb. 2018 Preg.~2 & Capacidad ilimitada & Greedy de monedas: qué técnica es y cuándo falla.\\
Ex. feb. 2018 Preg.~4 & Teoría & Qué resuelve alineamiento de secuencias y con qué técnica.\\
Ex. feb. 2018 Prob.~B & Último elemento & LIS con $G_i$ y $L_i$; pila de bandejas.\\
Ex. jul. 2018 Prob.~A & Grilla & Recolector con calles flechadas; esquinas intransitables.\\
}

\pdmapa{
2P 2018 Ej.~2 & Cantidad de aristas & Costo mínimo con $\le K$ aristas vía $\mathit{CMC}(i,v)$; algoritmo y orden.\\
Ex. dic. 2018 Ej.~3 & Dos secuencias & Ciclos mínimos para ejecutar dos programas en paralelo.\\
Ex. feb. 2019 Ej.~2 & Dos secuencias & Distancia de edición; lema de las últimas posiciones.\\
Ex. jul. 2019 Ej.~2 & DAG (conteo) & Largo mínimo de las historias y cuántas hay, en $\Oh(n)$.\\
2P 2019 Ej.~2 & Prefijo binario & Ladrón de casas no contiguas; valor y conjunto óptimo.\\
Ex. dic. 2019 Ej.~2 & Prefijo con estado & Neumáticos A/B por vuelta con costo de parada; $\OPT(i,X)$.\\
Ex. feb. 2020 Ej.~3 & Sufijo booleano & Dinosaurio que avanza 1 o salta 4: ¿hay secuencia ganadora?\\
Ex. jul. 2020 Ej.~3 & Capacidad ilimitada & Varillas con la mínima cantidad de soldaduras; $C(i,\ell)$.\\
2P 2020 Ej.~2 & Segmentación & Fragmentos con costo máximo menos mínimo y costo por corte.\\
Ex. dic. 2020 Ej.~3 & Capacidad ilimitada & Cambio con menos billetes; instancia donde greedy falla.\\
Ex. jul. 2021 Ej.~2 & Intervalos & Producto de una cadena de matrices.\\
2P 2021 Ej.~1 & Sufijo con salto & Tarifa por consumo o plana de 7 días; planificación óptima.\\
}

\pdmapa{
Ex. feb. 2022 Ej.~1 & Árbol & Independiente de peso máximo en árbol binario ($p$ y $q$); orden $\Oh(n)$.\\
Ex. feb. 2022 Ej.~3 & Grilla & Perforación minera de ganancia máxima; plan de perforación.\\
Ex. jul. 2022 Ej.~2 & Intervalos & Subsecuencia palíndroma más larga.\\
2P 2022 Ej.~1 & Prefijo con estado numérico & Paradas de F1 con desgaste de neumáticos; $\OPT(i,j)$.\\
Ex. feb. 2023 Ej.~3 & DAG & Camino más largo con orden topológico dado, en $\Oh(|V|+|E|)$.\\
Ex. jul. 2023 Ej.~2 & Cantidad de aristas & Routers: tiempo mínimo con a lo sumo $i$ saltos; camino.\\
2P 2023 Ej.~1 & Grilla & Sargento: costo mínimo hasta la fila superior; movimientos.\\
Ex. dic. 2023 Ej.~2 & Último elemento & Subsecuencia creciente más larga y reconstrucción.\\
Ex. feb. 2024 Ej.~2 & Segmentación & Concatenación de prefijos con máximas coincidencias.\\
2P 2024 Ej.~1 & Prefijo con estado & Cuchillas A y B con costo de cambio; solución óptima.\\
Ex. dic. 2024 Ej.~2 & Sufijo con salto & Ravioles y sorrentinos; planificación.\\
2P 2025 Ej.~3 & Sufijo con salto & Pedidos de soporte de uno o dos días.\\
Ex. feb. 2026 Ej.~3 & Segmentación & Segmentación mínima en palíndromos; posiciones de inicio.\\
Ex. jul. 2026 Ej.~2 & DAG & Tiempos de inicio con dependencias (camino crítico) en $\Oh(n+m)$.\\
}
